\documentclass[11pt]{article}
\usepackage[T1]{fontenc}
\usepackage[utf8]{inputenc}
\usepackage[margin=1in]{geometry}
\usepackage{amsmath,amssymb,amsthm,mathtools}
\usepackage{booktabs,array}
\usepackage[expansion=false]{microtype}
\usepackage{xcolor}
\usepackage[colorlinks=true,linkcolor=blue,citecolor=blue,urlcolor=blue]{hyperref}
\hypersetup{pdftitle={Short sextic character families: extraction and signed reductions},
 pdfauthor={Drafted for Edward Baker}}
\allowdisplaybreaks[2]
\setlength{\emergencystretch}{3em}
\numberwithin{equation}{section}
\newtheorem{theorem}{Theorem}[section]
\newtheorem{proposition}[theorem]{Proposition}
\newtheorem{lemma}[theorem]{Lemma}
\newtheorem{corollary}[theorem]{Corollary}
\theoremstyle{definition}
\newtheorem{definition}[theorem]{Definition}
\newtheorem{assumption}[theorem]{Analytic input}
\theoremstyle{remark}
\newtheorem{remark}[theorem]{Remark}
\newcommand{\N}{\mathrm N}
\newcommand{\OO}{\mathcal O}
\newcommand{\muK}{\mu_K}
\newcommand{\ind}{\mathbf 1}
\newcommand{\rad}{\operatorname{rad}}
\newcommand{\Mfrak}{\mathfrak M}
\newcommand{\eps}{\varepsilon}
\newcommand{\Res}{\operatorname*{res}}
\newcommand{\Rea}{\operatorname{Re}}
\newcommand{\Q}{\mathbb Q}
\newcommand{\C}{\mathbb C}
\newcommand{\R}{\mathbb R}
\newcommand{\dd}{\,d}
\newcommand{\good}{\mathcal I_S}
\newcommand{\sfstar}{\sideset{}{^*}\sum}
\newcommand{\norm}[1]{\left\lVert #1\right\rVert}
\title{Short sextic character families:\\
 extraction, adaptive factorization, and signed remainders}
\author{Drafted for Edward Baker}
\date{8 October 2026}
\begin{document}
\maketitle
\begin{abstract}
We organize reductions for smoothed M\"obius sums twisted by sextic
residue characters over the Eisenstein field. A moment at one fixed
positive row exponent \(h\), with loss \(a\), conditionally yields the
scalar exponent \((1+a)/2+5h/12\); finite reuse of the same moment
removes the prime-deletion error without an initial zero-free strip.
An exact truncated-inverse factorization and zero-integral derivative
profiles lead to a sufficient adaptive tail moment. Its cutoff depends
on the combined conductor and deletion modulus, which is comparable
to the good radical of the row. The actual omitted sector has
arbitrarily rapid mean-square decay with the complete Schwartz row
weight. We also prove a signed gcd truncation and uniform complete-kernel
bounds controlling a larger transformed overlap sector.
A further exact divisor cancellation makes a growing nonunit-cofactor
sector negligible under the complete row weight, even though its opposite
branches can each be oversized. Unit-cofactor semiprimes survive this
cancellation; it gives no improved bound for the full moment.
A square-divisor projection and sparse use of the generic physical
operator bound the entire nonsquarefree tail sector by
\(D^{19/24+\varepsilon}\) at \(H=D^{2/5}\), with cutoff
\(\theta=11/20\). This is below the proposed \(D^{4/5}\) budget
and leaves an equivalent squarefree signed target.
Within that target, complete small-prime cofactors cancel odd rough cores.
A growing three-prime sector is negligible on a specified large-conductor
range; even rough cores and the low-conductor response remain.
Prime-selected and prime-by-prime factor-selected obstructions show
why separate absolute estimates cannot give a useful moment;
the generic sextic large sieve retains the same coherent-row power.
The arithmetic tail estimate remains open. No new zero-free region
or implication from quasi-RH to RH is established.
\end{abstract}
\tableofcontents

\section{Introduction and the remaining estimate}\label{sec:intro}

The objective is to turn cancellation in a short family of sextic
twists into a stronger bound for one fixed-character M\"obius sum.
The distinction between a reduction and a new arithmetic estimate is
essential. Repeated rows make extraction effective, but they also
create a coherent obstruction for coefficient classes that discard
the M\"obius signs.

For \(K=\Q(\sqrt{-3})\), fixed arithmetic data \((\nu,S)\), and a smooth
annular profile \(W\), the amplitudes and normalized moment are
\[
 A_{u,W}(D)=\sum_{n\in\good}\muK(n)\nu(n)\chi_n(u)W(\N n/D),
 \qquad
 \Mfrak_W(D,H)=\frac1D\sum_{u\ne0}
       \Phi(\N u/H)|A_{u,W}(D)|^2.
\]
All notation and analytic inputs are specified in
Section~\ref{sec:setting}. The desired new estimate is
\begin{equation}\label{eq:intro-moment}
 \Mfrak_W(D,D^h)\ll_{\eps,W}D^{h+a+\eps},
 \qquad a+\frac{5h}{6}<\frac34.
\end{equation}
The inequality on \(h,a\) is exactly the condition for the extracted
scalar exponent to improve the benchmark \(7/8\).
The benchmark is used only for comparison; no zero-free theorem
from a neighboring investigation is assumed in the extraction.

The useful results fit together as follows.
Theorem~\ref{thm:finite-extraction} extracts a scalar bound from one fixed
moment. The exact factorization and adaptive small-sector estimate
in Section~\ref{sec:factorization} then replace
\eqref{eq:intro-moment} by a moment of an actual convolution tail.
Profiles
\[
 W=(1+t\partial_t)V,\qquad V\in C_c^\infty((0,\infty)),
\]
make the principal correction vanish and retain every Mellin detector
away from \(s=1\). Signed gcd recombination and transformed overlap
bounds give parallel formulations in
Sections~\ref{sec:gcd} and~\ref{sec:overlap}.
Section~\ref{sec:packets} proves exact signed cancellation in a growing
family of nonunit-cofactor packets, with full Schwartz-weighted control.
It retains the unit-cofactor semiprimes and does not prove a smaller full
moment. The generic sieve and selected-factor obstructions identify the
cancellation still required.
Section~\ref{sec:squarefree-tail} proves a power bound for the whole
nonsquarefree sector at a larger cutoff buffer, leaving a squarefree
signed residual with pairwise-coprime factors. This component bound
does not improve the proved exponent of the full moment.

One concrete sufficient target is the adaptive tail moment at
\(h=4/5\), \(0\le a<1/12\), for every derivative profile and one
fixed cutoff buffer \(\theta>0\). It would give
\[
 |A_{1,W}(D)|\ll_{\eps,W}D^{5/6+a/2+\eps},
 \qquad \frac56+\frac a2<\frac78.
\]
The present results simplify this target but do not prove it.

The manuscript develops the useful conclusions of research
notes~\cite{N1,N2,N3,N4,N5,N6,N7,N8,N9,N10,N11,N12,N13}
and the signed cancellation and squarefree reductions of~\cite{N15,N17}.
The lengthy ray-group identities are isolated as explicit imported
inputs; their conventions and detailed derivations remain in the
cited notes. Routes that yield no additional power saving are
summarized in Section~\ref{sec:limits}.

\section{Arithmetic setting and analytic inputs}\label{sec:setting}

\subsection{Ideals, rows and zero extensions}

Let \(\OO=\mathbb Z[\omega]\), \(\omega=e^{2\pi i/3}\).
For an element \(u\), \(\N u=u\bar u\); for a nonzero integral
ideal \(n\), \(\N n=\#(\OO/n)\).
Fix a finite-order Hecke character \(\nu\) and a finite set \(S\)
containing the primes above \(6\) and its conductor.
The monoid \(\good\) consists of integral ideals prime to \(S\).
Since \(\OO\) is a principal ideal domain, ideal variables may be
represented by fixed primary generators when a residue-symbol
argument requires an element. Rows \(u\) always range over elements.

At a good prime ideal \(p\), the native sextic symbol
\(\chi_p(u)\in\mu_6\cup\{0\}\) is the sixth-power residue symbol,
zero if \(p\mid(u)\). Fix one common orientation and put
\[
 \chi_n(u)=\prod_{p\mid n}\chi_p(u)^{v_p(n)},\qquad
 \chi_1(u)=1.
\]
This is the multiplicative denominator extension, including
nonsquarefree \(n\).
For fixed \(u\ne0\),
\begin{equation}\label{eq:lambda}
 \lambda_u(n)=\nu(n)\chi_n(u)
\end{equation}
is completely multiplicative, including its zero values.
In particular
\begin{equation}\label{eq:power-mask}
 \chi_n(r^6)=\ind_{(n,r)=1},\qquad
 \chi_n(vr^6)=\chi_n(v)\ind_{(n,r)=1}.
\end{equation}
The first identity is a deletion mask, not a constant character.
All convolutions below take place on \(\good\).
The functions \(\muK,\tau,\omega\) denote ideal M\"obius, divisor
count and number of distinct prime factors.

Fix \(W\in C_c^\infty((0,\infty))\), with
\(\operatorname{supp}W\subset[c,C]\), \(0<c<C\).
We allow signed and complex profiles.
Fix a nonnegative radial Schwartz weight \(\Phi\), at least one on
\([0,1]\); the transformed formulation additionally uses a smooth
compactly supported radial Fourier transform \(\widehat\Phi\).
Such weights exist: take a nonnegative radial smooth function of
sufficiently small Fourier support, square the modulus of its inverse
transform, and normalize its positive lower bound on the unit norm ball.

Write \(A_u=A_{u,W}\) when the profile is fixed, and set
\begin{equation}\label{eq:moment}
 \Mfrak_W(D,H)=\frac1D\sum_{u\in\OO\setminus\{0\}}
    \Phi(\N u/H)|A_{u,W}(D)|^2.
\end{equation}
We take \(D\ge2\) and \(1\le H\le D\) unless stated otherwise.
The zero row would contribute nothing for large \(D\):
all nonunit columns vanish there and the unit is off the profile.
All implied constants may depend on the fixed arithmetic data and
the displayed profiles. Every epsilon estimate means \emph{for every
\(\eps>0\)}, with its own constant and sufficiently large threshold.
No uniformity in a moving character or profile is implicit.

\subsection{The precise imported facts}

Elementary counting gives
\[
 \#\{n:\N n\le X\}\ll X,\quad
 \#\{u:0<\N u\le X\}\ll X,\quad
 \sum_{\N n\le X}\tau(n)\ll X\log(2X),\quad
 \tau(n)\ll_\eps(\N n)^\eps.
\]
The fixed-field prime ideal theorem gives
\(\#\{p:X/2<\N p\le X,\ p\notin S\}\asymp X/\log X\).
The more precise remainder needed for the selected-factor obstruction
is stated separately in Section~\ref{sec:obstructions}.

We retain the following arithmetic and Poisson package from the native
sextic presentation and the source conventions audited in
\cite{N2,N3,N8,N11,OA5}.

\begin{assumption}[Native presentation and primitive completion]
\label{input:poisson}
For each \(u\ne0\), the ideal character \(\lambda_u\) is induced by a
primitive finite-order Hecke character of conductor norm \(Q_u\).
Let \(E_u\) be the additional squarefree deletion ideal, with support
disjoint from that conductor, and put \(L_u=Q_u\N E_u\).
At a good prime dividing \(u\), a valuation nonzero modulo six gives
conductor exponent one; a positive valuation zero modulo six gives
only a deletion prime. The fixed bad part is bounded. Consequently
\begin{equation}\label{eq:radical-modulus}
 \N\rad_S(u)\le L_u\ll_{\nu,S}\N\rad_S(u)\le \N u,
 \qquad Q_u\le L_u,
\end{equation}
where \(\rad_S(u)\) is the product of good primes dividing \((u)\).

Primitive lattice Poisson has normalized Gauss factors of modulus one.
For a divisor \(e\mid E_u\), a radial sum at norm scale \(X\)
has nonzero-frequency prefactor of magnitude
\(X/(\N e\sqrt{Q_u})\) and norm frequency scale
\(X/(\N e Q_u)\). Its zero frequency is the principal residue times
\(X\int W\).
For two distinct squarefree columns \(m_1,m_2\), their ratio character
is nonprincipal with good conductor
\(m_1m_2/(m_1,m_2)^2\), and the canceled shared primes remain as a
coprimality mask.
\end{assumption}

The proofs below derive all-scale bounds from this package; they do
not assume a new conductor-uniform M\"obius bound.
Exact Fourier constants in the involution use the same normalization
as the imported paired-Gauss identity in
Section~\ref{sec:overlap}. Fixed ideal-to-lattice and bad-prime constants
are included in the residue and transforms.

For the Mellin conclusion we use the standard meromorphic continuation
of \(L_K^S(s,\nu)\) and its Euler product in \(\Rea s>1\).
A fixed-character boundary \(A_{u,V}(t)=o(t)\) is used only for the
optional sharp-profile recovery theorem. The generic sieve and the
source's positive dual theorem are explicitly identified where used.
Their deep proofs are not part of this manuscript.

% Drafted for Edward Baker, 8 October 2026, with substantial LLM assistance.
% Model: GPT-6 (Codex), inherited configuration; effort is not exposed.
% Manuscript fragment: uses amsmath, amssymb, amsthm and standard theorem environments.
% Imported analytic facts are identified in the text; the short-family target is not proved.

\section{The physical family and conditional extraction}
\label{sec:physical-extraction}

Retain the setting of Section~\ref{sec:setting}, and write
\(A(D)=A_{1,W}(D)\) for the fixed profile.
The squarefree support comes from \(\muK\); it is not an extra
restriction on subsequently introduced convolution factors.

The sharp normalized moment is
\[
 \Mfrak_{\le}(D,H)=D^{-1}\sum_{0<\N u\le H}|A_u(D)|^2.
\]
If $\Phi$ is a fixed nonnegative radial Schwartz weight, at least one
on $[0,1]$, its smoothed analogue
$D^{-1}\sum_{u\ne0}\Phi(\N u/H)|A_u(D)|^2$ dominates this sharp moment.
Thus a bound for the smoothed moment with the same exponent suffices in
the following theorem. No converse comparison at one fixed short-family
exponent is assumed.

\begin{theorem}[Finite extraction from one fixed short-family moment]
\label{thm:finite-extraction}
Fix $0<h\le1$ and $a\ge0$. Suppose, for every $\eps>0$ and all
sufficiently large real $D$, that
\begin{equation}\label{eq:short-moment-hypothesis}
 \Mfrak_{\le}(D,D^h)\ll_{\nu,S,W,h,a,\eps}
 D^{h+a+\eps}.
\end{equation}
Then, with
\[
 \beta(h,a)=\frac{1+a}{2}+\frac{5h}{12},
\]
one has $A(D)\ll_{\nu,S,W,h,a,\eps}
D^{\beta(h,a)+\eps}$ for every $\eps>0$.
The hypothesis is required at one fixed $h$, for one fixed character
and profile. An initial zero-free strip, a supremum over smaller column
scales, and conductor-uniform estimates are not needed for this deduction.
\end{theorem}

\begin{proof}
Put $B_p(L)=A_{p^6}(L)$ and $q=\N p$. Splitting squarefree ideals
according to divisibility by $p$, using \eqref{eq:power-mask},
gives the exact identities
\begin{equation}\label{eq:prime-completion}
 A(L)=B_p(L)-\nu(p)B_p(L/q),\qquad
 B_p(L)=\sum_{j\ge0}\nu(p)^jA(L/q^j).
\end{equation}
The second sum is finite because $A(L)=0$ for $L<1/C$.
Let $P(t)$ mean $A(L)\ll_\eps L^{t+\eps}$ for every
$\eps>0$ and $L\ge1$. Its constant may be enlarged to include
the bounded nonempty scales below one. For $t>0$, the geometric sum in
\eqref{eq:prime-completion} proves, uniformly for $q\ge2$,
\begin{equation}\label{eq:prime-mask-error}
 P(t)\quad\Longrightarrow\quad
 |B_p(D)-A(D)|\ll_{t,\eps}(D/q)^{t+\eps}.
\end{equation}

Set $Q=D^{h/6}$. We use the prime ideal theorem for the fixed field,
for example the fixed-field specialization of \cite{DKN}.
The number $J(D)$ of good prime ideals with $Q/2<\N p\le Q$ satisfies
$J(D)\asymp Q/\log Q$. Their sixth powers are distinct physical rows
of norm at most $D^h$. Hence \eqref{eq:short-moment-hypothesis} implies
\[
 J(D)^{-1}\sum_{Q/2<\N p\le Q}|B_p(D)|^2
 \ll_\eps D^{1+a+5h/6+\eps}
 =D^{2\beta(h,a)+\eps}.
\]
The logarithm has been absorbed by requesting a smaller preliminary
$\eps$. Applying \eqref{eq:prime-mask-error}, squaring and
averaging therefore gives
\[
 P(t)\Longrightarrow P\bigl(\max\{\beta(h,a),(1-h/6)t\}\bigr).
\]
Ideal counting supplies $P(1)$. If $\beta(h,a)\ge1$ it already proves
the assertion. Otherwise write $c_h=1-h/6\in(0,1)$. After $j$ uses of
the same hypothesis the exponent is
$t_j=\max\{\beta(h,a),c_h^j\}$; this follows immediately by induction.
A finite $j$ has $c_h^j\le\beta(h,a)$ and proves the desired exponent.
The arbitrary-$\eps$ quantifier absorbs each of the finitely many
losses; no uniformity in the iteration number is required.
\end{proof}

Thus a single pair improves $7/8$ precisely when
\begin{equation}\label{eq:improvement-budget}
 a+5h/6<3/4.
\end{equation}
For example, the lossless floors at $h=8/9,4/5,2/3$ are respectively
$47/54,5/6,7/9$. The corresponding strict losses allowed to improve
$7/8$ are $a<1/108,1/12,7/36$.

\subsection{Zero detection and the scope of replication}

For $0<\beta\le1$, the proof also gives the exact equivalence
\begin{equation}\label{eq:prime-row-equivalence}
 P(\beta)\quad\Longleftrightarrow\quad
 J(D)^{-1}\sum_{Q/2<\N p\le Q}|B_p(D)|^2
 \ll_\eps D^{2\beta+\eps}.
\end{equation}
The forward implication uses the geometric completion; the reverse
implication repeats the finite argument with floor $\beta$.
The positivity $\beta>0$ is essential here. This repeated-row estimate
is consequently a reformulation of its fixed-character target, rather
than an independent cancellation estimate.

To obtain a zero-free consequence one must assume the moment for a
class of profiles with no common Mellin zero in the region in question.
Indeed, with $\widehat W(s)=\int_0^\infty W(t)t^{s-1}\,dt$,
\begin{equation}\label{eq:extraction-mellin}
 \int_0^\infty A(D)D^{-s}\,\frac{dD}{D}
 =\frac{\widehat W(s)}{L_K^S(s,\nu)}\qquad(\Rea s>1).
\end{equation}
This follows by absolute convergence and the ideal M\"obius Euler
product. A bound $P(\beta)$ makes the left side holomorphic in
$\Rea s>\beta$. Multiplying by $L_K^S$ and continuing the identity
excludes any zero $\rho\ne1$ with $\Rea\rho>\beta$ for which
$\widehat W(\rho)\ne0$. The possible principal pole at $s=1$ is
handled meromorphically, and omitted Euler factors are nonzero for
$\Rea s>0$. Each zero and its detecting profile are fixed before the
constant and sufficiently large threshold are chosen.

Lemma~\ref{sfac:lem:profiles} below supplies a sufficient derivative
profile class with common Mellin zero only at \(s=1\).

\begin{remark}[Limits of replication alone]
Put $B_r(D)=A_{r^6}(D)$ for good integral ideals $r$.
The exact composite completion is
$B_r(D)=\sum_{\rad(d)\mid\rad(r)}\nu(d)A(D/\N d)$,
with arbitrary prime powers allowed. Its inverse is the finite
squarefree-divisor sum
$A(D)=\sum_{d\mid\rad(r)}\muK(d)\nu(d)B_r(D/\N d)$.
For every fixed $\beta>0$, the product
$\prod_{p\mid r}(1-(\N p)^{-\beta-\delta})^{-1}$ is
$O_{\beta,\delta,\eta}((\N r)^\eta)$ for every $\delta,\eta>0$.
This proves that averaging all $\N r\le D^{h/6}$ is equivalent, up to
arbitrary epsilon losses, to the same $P(\beta)$ target; the reverse
implication retains the large-prime subsum. Composite rows improve the
prime count only logarithmically. A fixed small prime divisor of $r$
prevents replacing their masks by the false uniform error
$O((D/\N r)^t)$.
For arbitrary weights on $R$ selected sixth-power rows,
$|\sum_r w_r|^2/\sum_r|w_r|^2\le R\ll D^{h/6}$.
For $d\mid6$, the inherited order-dividing-$d$ subfamily has
$h_d=dh/6$, $a_d=a+h-h_d$, and therefore the same extraction floor
$(1+a_d)/2+(d-1)h_d/(2d)=\beta(h,a)$.
These are limits of extraction from the same unweighted moment, not
no-go theorems for new signed correlation estimates. Fixed nontrivial
cores replicate twisted targets; growing cores require conductor-uniform
input. Detailed proofs and qualifications are in \cite{N5,N6}.
\end{remark}


% Integrable manuscript section; not a standalone document.
% Drafted for Edward Baker, 8 October 2026, with substantial LLM assistance.
% Model: GPT-6 (Codex), inherited configuration; exact variant and effort
% label are not exposed.  This is a same-model internal derivation.
% Requires amsmath, amssymb, amsthm and theorem/lemma/proposition/corollary
% environments.  All labels are prefixed sfac:.

\section{A zero-integral adaptive factorization}
\label{sfac:sec:adaptive}\label{sec:factorization}

Retain \(A_{u,W},\lambda_u,\Mfrak_W\) from
Section~\ref{sec:setting}. All convolutions are on \(\good\).
The symbols \(\delta\) and \(\mathbf1\) denote the convolution identity
and the constant-one function. Complete multiplicativity includes
every deletion zero. The arguments in this section develop
\cite{N8,N10,N11}.

\subsection{An exact truncated inverse}

\begin{lemma}[Truncated Dirichlet inverse]
\label{sfac:lem:inverse}
For $z\ge1$ set $m_z=\mu_K\mathbf1_{\N n\le z}$ and
$r_z=\delta-\mathbf1*m_z$.  For every integer $J\ge1$,
\begin{equation}
 \mu_K=\sum_{j=1}^J(-1)^{j-1}\binom Jj
            m_z^{*j}*\mathbf1^{*(j-1)}+\mu_K*r_z^{*J}.
 \label{sfac:eq:inverse}
\end{equation}
The first sum equals $\mu_K(n)$ at every $\N n\le z^J$, including
the endpoint.
\end{lemma}

\begin{proof}
The equality $\mathbf1*\mu_K=\delta$ shows that $r_z(n)=0$ when
$\N n\le z$.  Since $\delta-r_z=\mathbf1*m_z$, the binomial identity
\[
 \delta-r_z^{*J}
  =\sum_{j=1}^J(-1)^{j-1}\binom Jj(\mathbf1*m_z)^{*j}
\]
gives \eqref{sfac:eq:inverse} after convolution with $\mu_K$.
Every nonzero term of $r_z^{*J}$ is a product of $J$ ideals whose
norms are strictly greater than $z$.  Convolution with $\mu_K$ does
not decrease the norm.  Its remainder is therefore supported strictly
beyond $z^J$.
\end{proof}

Multiplication of coefficients by $\lambda_u$ commutes with convolution.
This assertion remains true when factors overlap: zero extension is part
of complete multiplicativity.  No squarefree or pairwise-coprimality
restriction may be added to the individual factors of
\eqref{sfac:eq:inverse}.  Their signed recombination supplies the
required nonsquarefree zeros.

Take $J=2$, $z=(CD)^{1/2}$, and $D$ sufficiently large that $cD>z$.
Then $m_z$ vanishes on the profile, and, with $c_z=m_z*m_z$,
\begin{align}
 S_u(X)&=\sum_{m\in\good}\lambda_u(m)W(\N m/X),\label{sfac:eq:S}\\
 A_{u,W}(D)&=-\sum_d c_z(d)\lambda_u(d)S_u(D/\N d).
 \label{sfac:eq:factor}
\end{align}
All these sums are finite; $c_z(d)$ is supported on $\N d\le CD$ and
$|c_z(d)|\le\tau(d)$.  It is cube-free supported.  For $\N d\le z$,
$c_z(d)=(\mu_K*\mu_K)(d)$, whose local coefficients are $-2$ at
$p$, $1$ at $p^2$, and $0$ at $p^j$ for $j\ge3$.
The truncation $z^2\ge CD$ covers the entire profile, rather than
only its lower portion.

\subsection{A sufficient zero-integral profile class}

\begin{lemma}[Profiles with vanishing principal density]
\label{sfac:lem:profiles}
If $V\in C_c^\infty((0,\infty))$ and
\begin{equation}
 W(t)=(1+t\partial_t)V(t),
 \label{sfac:eq:derivative-profile}
\end{equation}
then $I_W:=\int_0^\infty W(t)\,dt=0$ and
\begin{equation}
 \widehat W(s)=(1-s)\widehat V(s),\qquad
 \widehat V(s)=\int_0^\infty V(t)t^s\,\frac{dt}{t}.
 \label{sfac:eq:profile-Mellin}
\end{equation}
The common Mellin zero of this profile class is exactly $s=1$.
\end{lemma}

\begin{proof}
Both identities follow by integration by parts, with vanishing endpoint
terms.  Given $\rho\ne1$, choose a nonzero nonnegative smooth annular
$\phi$ and put $V(t)=t^{-\rho}\phi(t)$.  Then
$\widehat W(\rho)=(1-\rho)\int\phi(t)\,dt/t\ne0$.
\end{proof}

Consequently these profiles suffice in the scalar Mellin detector
$\widehat W(s)/L_K^S(s,\nu)$ for every potential zero other than
$s=1$, which retains its classical separate treatment.  The finite
prime-mask extraction uses the same profile at each scale and applies
unchanged to signed or complex profiles.  Thus restricting the analytic
target to \eqref{sfac:eq:derivative-profile} removes no relevant detector;
it does not itself prove a short-family estimate.


\subsection{Sharp moment recovery for arbitrary profiles}
\label{sec:profile-recovery}

The derivative class is already sufficient for Mellin detection.
It also recovers every sharp profile moment in the useful range,
subject to a classical fixed-character boundary.

\begin{proposition}[Recovery with a fixed-row boundary]
\label{prop:sharp-recovery}
Let \(0<h<1\), \(a\ge0\), \(h+a<1\), and
\(W=(1+t\partial_t)V\), \(V\in C_c^\infty((0,\infty))\).
Assume the fixed-character M\"obius boundary
\begin{equation}\label{eq:fixed-boundary}
 A_{u,V}(t)/t\longrightarrow0\quad(t\longrightarrow\infty)
 \quad\text{for each fixed }u\ne0.
\end{equation}
If
\[
 \sum_{0<\N u\le t^h}|A_{u,W}(t)|^2
       \ll_{\eps,W}t^{1+h+a+\eps},
\]
then the same sharp moment bound holds with \(V\) in place of \(W\).
Thus the sharp all-profile and derivative-profile moment assertions
are equivalent for \(h+a<1\), subject to \eqref{eq:fixed-boundary}.
\end{proposition}
\begin{proof}
Termwise differentiation of the locally finite smooth sum gives
\(A_{u,W}(t)=(1-t\partial_t)A_{u,V}(t)\).
Since \((A_{u,V}(t)/t)'=-A_{u,W}(t)/t^2\),
the boundary assumption yields
\begin{equation}\label{eq:profile-integral-recovery}
 A_{u,V}(D)=D\int_D^\infty A_{u,W}(t)\,\frac{dt}{t^2}.
\end{equation}
For the finite fixed row set \(0<\N u\le D^h\), Minkowski and
enlargement to \(0<\N u\le t^h\), \(t\ge D\), give
\[
 \left(\sum_{0<\N u\le D^h}|A_{u,V}(D)|^2\right)^{1/2}
 \ll D\int_D^\infty t^{(1+h+a+\eps)/2-2}\,dt
 =\frac{D^{(1+h+a+\eps)/2}}
        {1-(1+h+a+\eps)/2}.
\]
Choose a smaller preliminary epsilon so that \(h+a+\eps<1\).
Squaring proves the result. The reverse implication is restriction
of the all-profile assertion.
\end{proof}

The ordinary fixed-character M\"obius prime number theorem supplies
\eqref{eq:fixed-boundary}; here it is an explicitly retained analytic
input, used only for this optional recovery. Its constants need not
be conductor-uniform: the boundary is used for finitely many fixed
rows at each fixed \(D\). This argument recovers a sharp moment,
not an arbitrary Schwartz-weighted all-row moment.
A bound for the latter on derivative profiles supplies the sharp
hypothesis because \(\Phi\ge1\) on the inner ball.

\subsection{Native arithmetic inputs and primitive completion}

Use Analytic input~\ref{input:poisson}, and extend \(\lambda_u\)
by zero at \(S\). For the primitive inducing character \(\psi_u\),
\begin{equation}\label{sfac:eq:primitive-presentation}
 \lambda_u(n)=\psi_u(n)\ind_{(n,E_u)=1},\qquad
 Q_u=\\N\mathfrak q_u,\quad L_u=Q_u\N E_u.
\end{equation}
The radical comparison is \eqref{eq:radical-modulus}.
We retain the principal frequency through its exact residue and derive
the remaining all-scale bound from primitive Poisson.

\begin{lemma}[All-scale masked completion]
\label{sfac:lem:completion}
Let $\kappa_u$ be the residue at $s=1$ of the Dirichlet series of
$\lambda_u$.  Then
\begin{equation}
 S_u(X)=\kappa_u XI_W+\mathcal V_u(X),\qquad
 |\mathcal V_u(X)|\ll_{A,W}
      \tau(E_u)\sqrt{Q_u}(L_u/X)^A
      \quad(X>0,\ A\ge0).
 \label{sfac:eq:completion}
\end{equation}
Here $\kappa_u=0$ unless $\psi_u$ is principal.  In that case
\begin{equation}
 \kappa_u=\frac{\pi}{3\sqrt3}
       \prod_{p\in S\,\text{or}\,p\mid u}(1-(\N p)^{-1}).
 \label{sfac:eq:residue}
\end{equation}
The remainder $\mathcal V_u$ is exactly the full nonzero-frequency
Poisson remainder, without a row or frequency truncation.
\end{lemma}

\begin{proof}
For a primitive character of conductor norm $Q$, primitive Poisson after
removal of its zero frequency has prefactor of size $\sqrt Q\,t$,
where $t=X/Q$.  Its remaining lattice frequencies have radial Schwartz
weights at scale a fixed constant times $t$.  Lattice counting gives
\[
 \sum_{k\ne0}|\mathcal F_W(t\N k)|\ll_W t^{-1}\quad(t>0),
 \qquad
 \sum_{k\ne0}|\mathcal F_W(t\N k)|\ll_{A,W}t^{-A-1}
       \quad(t\ge1).
\]
The fixed field and Fourier-normalization constants are absorbed here.
Hence the primitive remainder is $O_W(\sqrt Q)$ at every scale and
$O_{A,W}(\sqrt Q(Q/X)^A)$ when $X\ge Q$.  For $X<Q$ the latter
bound follows from the former because $(Q/X)^A\ge1$.

Expand the mask in \eqref{sfac:eq:primitive-presentation}.  Each
$e\mid E_u$ multiplies the primitive sum by $\mu_K(e)\psi_u(e)$
and replaces $X$ by $X/\N e$.  Its remainder is bounded by
$\sqrt{Q_u}(\N e Q_u/X)^A$, and $\N e Q_u\le L_u$.
Summing the $\tau(E_u)$ divisors proves the error estimate.
The zero frequency occurs precisely for the principal primitive
character.  Its unmasked coefficient is
$\operatorname{res}_{s=1}\zeta_K(s)=\pi/(3\sqrt3)$; the divisor
expansion contributes the surviving Euler factors, giving
\eqref{sfac:eq:residue}.
\end{proof}

\begin{lemma}[Radical-weighted row mass]
\label{sfac:lem:mass}
For $H\ge1$ and every $\delta>0$,
\begin{equation}
 \sum_{u\ne0}\Phi(\N u/H)
       \frac{\tau(E_u)^2Q_u}{L_u^2}\ll_\delta H^\delta.
 \label{sfac:eq:mass}
\end{equation}
Also $\#\{0<\N u\le H:L_u\le R\}\ll_\delta RH^\delta$ for $R\ge1$.
\end{lemma}

\begin{proof}
The fixed bad part costs a constant, and
$\tau(E_u)^2Q_u/L_u^2\ll
4^{\omega(\operatorname{rad}_S(u))}/\N\operatorname{rad}_S(u)$.
For every $\delta>0$ the ideal Euler product
\begin{align}
 &\sum_{\mathfrak n}
 \frac{4^{\omega(\operatorname{rad}_S(\mathfrak n))}}
      {\N\operatorname{rad}_S(\mathfrak n)(\N\mathfrak n)^\delta}
 \notag\\[-2pt]
 &\hspace{1cm}=
 \prod_{p\notin S}
 \left(1+\frac{4(\N p)^{-1-\delta}}{1-(\N p)^{-\delta}}\right)
 \prod_{p\in S}(1-(\N p)^{-\delta})^{-1}
 <\infty\label{sfac:eq:Euler-mass}
\end{align}
converges absolutely.  The large-prime terms are $O((\N p)^{-1-\delta})$.
There are six element generators per ideal, and
$\Phi(\N u/H)\ll_\delta H^\delta(\N u)^{-\delta}$, proving
\eqref{sfac:eq:mass} with the complete Schwartz weight.

For the counting assertion, the same Euler product gives
\[
 \sum_{0<\N u\le H}\frac1{\\N\operatorname{rad}_S(u)}
       \ll_\delta H^\delta.
\]
On the counted rows $\N\operatorname{rad}_S(u)\le L_u\le R$.
This proves the count; the smaller mass in \eqref{sfac:eq:mass} alone
would not prove it on principal rows.
\end{proof}

\subsection{The adaptive tail theorem}

\begin{theorem}[Negligible adaptive small sector]
\label{sfac:thm:adaptive}
Fix $0<\theta<1$ and let $Y_u=D^{1-\theta}/L_u$.  Define
\begin{align}
 I_u&=-\sum_{\N d\le Y_u}c_z(d)\lambda_u(d)S_u(D/\N d),
 &T_u&=-\sum_{\N d>Y_u}c_z(d)\lambda_u(d)S_u(D/\N d),\notag\\
 P_u&=-\kappa_uDI_W
       \sum_{\N d\le Y_u}\frac{c_z(d)\lambda_u(d)}{\N d}.
 \label{sfac:eq:adaptive-sums}
\end{align}
The short sums are empty if $Y_u<1$.  For every $A\ge0$, $\delta>0$,
\begin{equation}
 D^{-1}\sum_{u\ne0}\Phi(\N u/H)|I_u-P_u|^2
  \ll_{A,\delta}D^{1-2\theta-2A\theta}\log^2(2D)H^\delta.
 \label{sfac:eq:adaptive-energy}
\end{equation}
For $1\le H\le D$ this is $O_N(D^{-N})$ for every requested $N$,
choosing $A$ after $N$ and $\theta$.

If $W$ has the form \eqref{sfac:eq:derivative-profile}, then $P_u=0$.
For every fixed $0<h<1$ and $a\ge0$, the full moment estimate
\begin{equation}
 \mathfrak M_W(D,D^h)\ll_\varepsilon D^{h+a+\varepsilon}
 \label{sfac:eq:full-target}
\end{equation}
is therefore equivalent to the tail estimate
\begin{equation}
 D^{-1}\sum_{u\ne0}\Phi(\N u/D^h)|T_u|^2
       \ll_\varepsilon D^{h+a+\varepsilon}.
 \label{sfac:eq:tail-target}
\end{equation}
The tail estimate remains an analytic hypothesis.
\end{theorem}

\begin{proof}
Ideal counting gives $\sum_{\N d\le Y}|c_z(d)|\ll Y\log(2Y)$ for
$Y\ge1$.  For $\N d\le Y_u$, \eqref{sfac:eq:completion} has
$L_u\N d/D\le D^{-\theta}$.  As $Y_u\le D$, this gives
\[
 |I_u-P_u|\ll_{A,W}
 D^{1-\theta}L_u^{-1}\log(2D)
       \tau(E_u)\sqrt{Q_u}D^{-A\theta}.
\]
Square, divide by $D$, and apply Lemma~\ref{sfac:lem:mass}.
This proves \eqref{sfac:eq:adaptive-energy}.  The exact factorization
is $A_{u,W}=(I_u-P_u)+(T_u+P_u)$.  For zero-integral profiles $P_u=0$,
and the triangle inequality in the weighted $\ell^2$ norm proves the
equivalence of \eqref{sfac:eq:full-target} and \eqref{sfac:eq:tail-target}.
\end{proof}

The cutoff is legitimate because the factorization and completion are
pointwise in $u$.  Squaring it produces a row-dependent arithmetic
selector; it cannot be inserted into a previously completed row kernel
whose estimates require a row-independent selector.

\begin{corollary}[Uniform-cutoff centered estimate]
\label{sfac:cor:uniform}
Instead split \eqref{sfac:eq:factor} at a fixed $1\le Y\ll D$, and
write the resulting sums as $I_{u,Y},T_{u,Y}$, with
$P_{u,Y}=-\kappa_uDI_W\sum_{\N d\le Y}c_z(d)\lambda_u(d)/\N d$.
For $1\le H\le D$ and every fixed $A\ge0$,
\begin{equation}
 D^{-1}\sum_{u\ne0}\Phi(\N u/H)|I_{u,Y}-P_{u,Y}|^2
 \ll_{A,\varepsilon}D^\varepsilon
       \frac{H^2Y^2}{D}\left(\frac{HY}{D}\right)^{2A}.
 \label{sfac:eq:uniform-energy}
\end{equation}
For $H=D^h$, $Y=D^{1-h-\eta}$, and fixed $0<\eta<1-h$, this
has arbitrary power decay.  Separately, its $A=0$ branch is affordable
at $D^{h+a+\varepsilon}$ whenever $Y=D^y$ and $2y+h-1\le a$.
For a general profile the residual is $T_{u,Y}+P_{u,Y}$, with the
principal correction retained inside the same square.
\end{corollary}

\begin{proof}
Equations \eqref{eq:radical-modulus} and
\eqref{sfac:eq:completion}, with the divisor bound, imply
$|\mathcal V_u(X)|\ll_{A,\varepsilon}(\N u)^{1/2+\varepsilon}(\N u/X)^A$.
Summing $|c_z|$ below $Y$ gives
$|I_{u,Y}-P_{u,Y}|\ll Y^{1+\varepsilon}
(\N u)^{1/2+\varepsilon}(\N u Y/D)^A$.
Every higher moment of the fixed Schwartz row weight is finite;
lattice counting therefore proves \eqref{sfac:eq:uniform-energy},
after reallocating epsilon.  At the buffered cutoff its prefactor is
$D^{1-2\eta}$ and its decay factor is $D^{-2A\eta}$.  The unbuffered
claim follows directly by taking $A=0$.
\end{proof}

For general profiles the same exact argument gives the residual
$T_u+P_u$.  In terms of the full Poisson remainder it has the useful
identity
\begin{equation}
 T_u+P_u=-\kappa_uDI_W M_z(u)^2
       -\sum_{\N d>Y_u}c_z(d)\lambda_u(d)\mathcal V_u(D/\N d),
 \qquad
 M_z(u)=\sum_{\N a\le z}\frac{\mu_K(a)\lambda_u(a)}{\N a}.
 \label{sfac:eq:main-square}
\end{equation}
Indeed the zero frequencies of the short and long sums combine into
$-\kappa_uDI_W\sum_d c_z(d)\lambda_u(d)/\N d$, and finite convolution
makes this last sum $M_z(u)^2$.  It is an ordinary square, not
$|M_z(u)|^2$.  Neither term of \eqref{sfac:eq:main-square} is separately
estimated by the reduction.

On the tail, write $d=ab$ with $\N a,\N b\le\sqrt{CD}$.  Its exact support
implies
\begin{equation}
 \N m<CD^\theta L_u,\qquad
 \N a,\N b>D^{1/2-\theta}/(\sqrt C L_u).
 \label{sfac:eq:tail-support}
\end{equation}
If $u=\eta vr^6$, with $v$ sixth-power-free and $\eta$ a unit, then
$L_u\ll \N\operatorname{rad}(vr)\le \N v\,\N r
=(\N u)^{1/6}(\N v)^{5/6}$.  On coherent inner rows $u=r^6$, $\N u\le D^h$,
this gives
\begin{equation}
 Y_u\gg D^{1-\theta-h/6},\qquad
 \N m\ll D^{\theta+h/6},\qquad
 \N a,\N b\gg D^{1/2-\theta-h/6}.
 \label{sfac:eq:coherent-support}
\end{equation}
For a complete Schwartz moment these bounds are used pointwise or on
row shells; $\N u\le D^h$ is not a global replacement for the row variable.

The following proposed row exponents give these support ranges on
coherent inner rows.  Constants depending on the fixed annulus and
arithmetic data are suppressed.
\begin{center}
\begin{tabular}{@{}cccc@{}}
\toprule
$h$ & lower bound for $Y_u$ & upper bound for $\N m$
    & lower bound for each $\N a,\N b$\\
\midrule
$2/5$ & $D^{14/15-\theta}$ & $D^{1/15+\theta}$
      & $D^{13/30-\theta}$\\
$4/5$ & $D^{13/15-\theta}$ & $D^{2/15+\theta}$
      & $D^{11/30-\theta}$\\
$8/9$ & $D^{23/27-\theta}$ & $D^{4/27+\theta}$
      & $D^{19/54-\theta}$\\
\bottomrule
\end{tabular}
\end{center}
These are support statements; the remaining signed sums are not bounded
by the completion argument.

\section{Exact signed cancellation in divisor packets}
\label{sec:packets}

The adaptive tail admits a further exact cancellation, developed in
\cite{N15}. It removes a particular family of total products; it does
not improve the proved bound for the complete tail moment. We first
recombine all factor tuples with the same total ideal, preserving
their signs, phases and deletion zeros.

\subsection{Complete recombination by total ideal}

For $z\ge1$ and $Y>0$, put
\begin{equation}\label{eq:packet-coefficient}
 t_{z,Y}(n)=-\sum_{\substack{d\mid n\\\N d>Y}}c_z(d).
\end{equation}
Complete multiplicativity in \eqref{eq:lambda} gives the finite identity
\begin{equation}\label{eq:packet-recombination}
 T_{u,W}(D)=\sum_{n\in\good}
   t_{z,Y_u}(n)\lambda_u(n)W(\N n/D),
 \qquad Y_u=D^{1-\theta}/L_u.
\end{equation}
Thus $t_{z,Y_u}(n)$ includes every original tuple $n=abm$,
with $\N a,\N b\le z$ and $\N(ab)>Y_u$.
For the adaptive factorization, $z=\sqrt{CD}$ and sufficiently large
$D$ give $z<\N n\le z^2$ on the profile. Thus
Lemma~\ref{sfac:lem:inverse} with $J=2$
also gives
\begin{equation}\label{eq:packet-inverse-coefficient}
 t_{z,Y}(n)=\muK(n)+\sum_{\substack{d\mid n\\\N d\le Y}}c_z(d).
\end{equation}
Indeed $m_z(n)=0$ there and $(\mathbf1*c_z)(n)=-\muK(n)$.
No coprimality or squarefreeness conditions have been added to the
original factors in either identity.

\subsection{A pointwise cancellation theorem}

For a nonunit ideal $b$, write
$P^-(b)=\min_{p\mid b}\N p$.

\begin{theorem}[Nonunit cofactor packet]
\label{thm:divisor-packet}
Let $b\ne1$ be any good ideal, $B=\N b$, and let $q,r$ be
distinct good prime ideals with $(b,qr)=1$. Suppose
\begin{align}
 &\N q,\N r\le z,\qquad B\N q\le z,
       \qquad \N r\,P^-(b)>z,\label{eq:packet-factor-conditions}\\
 &\max\{B\N q,B\N r\}\le Y<\N(qr).
       \label{eq:packet-tail-conditions}
\end{align}
Then
\begin{equation}\label{eq:packet-zero}
 t_{z,Y}(bqr)=0.
\end{equation}
The statement allows nonsquarefree $b$ and holds with every complex
completely multiplicative twist and its deletion zeros.
\end{theorem}

\begin{proof}
A divisor of $bqr$ missing either $q$ or $r$ has norm at most
$\max\{B\N q,B\N r\}$ and is excluded from
\eqref{eq:packet-coefficient}. The surviving divisors are exactly
$qre$, $e\mid b$, since $\N(qr)>Y$.
Also $Y\ge B\N r\ge P^-(b)\N r>z$, so $q,r$ cannot occur
together in one truncated factor.

No prime of $b$ can share the factor containing $r$, by
$\N r\,P^-(b)>z$. Every divisor $e\mid b$ can share the
factor containing $q$, by $\N(qe)\le B\N q\le z$.
The only possible ordered factor pairs for $qre$ are consequently
$(qe,r)$ and $(r,qe)$, and
\[
 c_z(qre)=2\muK(e).
\]
For nonsquarefree $e$, both sides vanish. Summation gives
\[
 t_{z,Y}(bqr)=-2\sum_{e\mid b}\muK(e)=0,
\]
because $b\ne1$. Every tuple carries the same value
$\lambda_u(bqr)W(\N(bqr)/D)$, including when that value is zero.
The cancellation therefore precedes squaring and survives all
complex phases and masks.
\end{proof}

For a prime cofactor $b=p$, the opposite branches are particularly
simple. The divisor $d=qr$, with free factor $m=p$, has total
coefficient $-2$; the divisor $d=pqr$, with $m=1$, has total
coefficient $+2$. Each coefficient includes both ordered truncated
factor pairs. Their contributions cancel at the same total ideal.

\subsection{A growing sector under the complete row weight}

\begin{corollary}[Negligible full weighted packet sector]
\label{cor:packet-energy}
Fix $h=2/5$, $\theta=1/40$, and $0<\delta<1/40$.
Let $\mathcal P_\delta(D)$ be the row-independent set of good ideals
$n$ with $cD\le\N n\le CD$ admitting a representation $n=bqr$
such that
\begin{equation}\label{eq:packet-growing-set}
 1<B\le D^{1/40-\delta},\qquad B\N q\le z,
 \qquad z/7<\N r\le z,
 \qquad (b,qr)=1,
\end{equation}
where $q,r$ are distinct good primes and $B=\N b$.
Each total ideal occurs once in the set, even if its representation
is not unique. Define its actual response vector by
\begin{equation}\label{eq:packet-response}
 \mathcal P_{u,W}(D)=
 \sum_{n\in\mathcal P_\delta(D)}
 t_{z,Y_u}(n)\lambda_u(n)W(\N n/D).
\end{equation}
For every fixed smooth annular $W$ and every requested $N>0$,
\begin{equation}\label{eq:packet-energy}
 \frac1D\sum_{u\ne0}\Phi(\N u/D^{2/5})
       |\mathcal P_{u,W}(D)|^2\ll_{N,W,\Phi,\delta}D^{-N}.
\end{equation}
No zero-integral condition on $W$ is needed for this sector estimate.
\end{corollary}

\begin{proof}
Every good prime ideal has norm at least seven, since $S$ contains
the primes above six. Thus \eqref{eq:packet-growing-set} implies
$\N r\,P^-(b)>z$. On the enlarged row range
$\N u\le D^{17/40}$, \eqref{eq:radical-modulus} gives
\begin{align*}
 Y_u&\gg D^{22/40},&
 \max\{B\N q,B\N r\}&\le Bz\ll D^{21/40-\delta},\\
 \N(qr)&=\N n/B\ge cD^{39/40+\delta},&
 Y_u&\le D^{39/40}.
\end{align*}
The positive exponent margins dominate all fixed constants for
sufficiently large $D$. Theorem~\ref{thm:divisor-packet} applies
pointwise, so $\mathcal P_{u,W}(D)=0$ on this entire enlarged ball.

For the farther rows, the uniform bound
\begin{equation}\label{eq:packet-uniform-bound}
 |\mathcal P_{u,W}(D)|
 \ll_W\sum_{\N n\le CD}\sum_{d\mid n}|c_z(d)|
 \ll_W D\log^2(2D)
\end{equation}
follows by majorizing with all ideal triples of product norm at most
$CD$. Indeed, ideal counting followed by two harmonic ideal sums
bounds that count by $D\log^2(2D)$.
For every $A>1$, lattice counting and Schwartz decay give
\[
 \sum_{\N u>D^{17/40}}\Phi(\N u/D^{2/5})
 \ll_{A,\Phi}D^{2/5}D^{-(A-1)/40}.
\]
Square \eqref{eq:packet-uniform-bound}, divide by $D$, and choose
$A$ after the requested $N$. This proves \eqref{eq:packet-energy}
with the complete row weight, rather than only an inner-ball sum.
\end{proof}

The only conductor input in this proof is $1\le L_u\ll_{\nu,S}\N u$.
The sector can be removed from the full response by the weighted
triangle inequality. For general profiles, the principal correction in
\eqref{sfac:eq:main-square} remains inside the residual $T_u+P_u$.
The individual factor branches of the packet cannot be discarded
separately.

\subsection{Large opposite branches and the surviving obstruction}

\begin{proposition}[Oversized branches with exact compensation]
\label{prop:packet-branches}
In the parameter setting of Corollary~\ref{cor:packet-energy}, let
$\nu=1$, fix a good prime cofactor $p$, and let $W\ne0$ be
a fixed smooth annular profile, possibly complex. There are fixed
narrow prime windows $\mathcal W_D$ selecting products $pqr$ in
$\mathcal P_\delta(D)$ such that the actual tail branches
\begin{equation}\label{eq:packet-branch-amplitudes}
 \mathcal B^\pm_{u,p}(D)=\pm2\lambda_u(p)
 \sum_{(q,r)\in\mathcal W_D}
 \lambda_u(qr)W(P\N(qr)/D),\qquad P=\N p,
\end{equation}
have normalized weighted energy, each separately,
\begin{equation}\label{eq:packet-branch-energy}
 \gg_{W,p,\Phi}\frac{DH^{1/6}}{\log^4D}
 =\frac{D^{16/15}}{\log^4D},\qquad H=D^{2/5}.
\end{equation}
The minus branch has $d=qr,m=p$, and the plus branch has
$d=pqr,m=1$. Both include the two ordered truncated factor pairs.
Their sum is exactly zero on every row. The lower bound concerns each
selected branch, not an unrestricted complex branch sum.
\end{proposition}

\begin{proof}
Write $P=\N p$ and choose $t\in(c,C)$ with $W(t)\ne0$.
Take
\[
 \max\{t/C,1/7\}<\beta<1,\qquad
 \alpha=t/(C\beta)<1.
\]
Choose the same sufficiently narrow windows around
$\N q=\alpha z/P$ and $\N r=\beta z$ for both branches.
Their centers satisfy $P\N q<z$, $\N r>z/7$, and
$P\N q\N r=tD$. Shrinking the windows preserves the strict
inequalities and ensures a positive lower bound for
$\Rea(e^{-i\arg W(t)}W)$ throughout their product range.
For large $D$, the primes are distinct and avoid $p$, and the fixed
cofactor satisfies $P\le D^{1/40-\delta}$.

The fixed-field prime ideal theorem supplies
$\asymp_{W,p}D/\log^2D$ pairs in these windows. Restrict rows
to $u=v^6$, $\N v\le H^{1/6}$, with $(v,p)=1$.
The primes $q,r\asymp D^{1/2}$ cannot divide $v$, so all original
masks equal one. There are $\asymp_p H^{1/6}$ distinct such
physical rows, accounting for the finite unit multiplicity.
Each selected branch amplitude has modulus
$\gg_{W,p}D/\log^2D$ there. Square, use $\Phi\ge1$ on the
inner ball, count these rows, and divide by $D$ to obtain
\eqref{eq:packet-branch-energy}.

For this fixed cofactor, $\N(qr)\ge cD/P>D^{39/40}\ge Y_u$
for sufficiently large $D$, uniformly in every row. Both selected
branches therefore belong to the actual tail on every row.
The window conditions give $c_z(qr)=2$ and $c_z(pqr)=-2$,
so their amplitudes are precisely \eqref{eq:packet-branch-amplitudes}
and cancel identically. Outside the enlarged row ball, the complete
coefficient $t_{z,Y_u}(pqr)$ may contain additional divisors; it is
the sum of the two selected branches, rather than that complete
coefficient, which vanishes on all rows. Prime ideal counting is
used for the branch lower bound, not for the algebraic cancellation
or the negligible-sector estimate.
\end{proof}

\begin{remark}[No improved full moment is proved]
\label{rem:packet-open-moment}
The unit-cofactor case is different. For balanced distinct
semiprimes $qr$ with $\N q,\N r\le z\le Y_u<\N(qr)$,
all proper divisors lie below the tail threshold, and
\[
 t_{z,Y_u}(qr)=-c_z(qr)=-2.
\]
Every decomposition of the same semiprime has already been included;
its compensation must involve different total ideals. The packet
theorem therefore does not bound the remaining signed tail.

At $h=a=2/5$, $\theta=1/40$, the proposed full tail target is
$D^{4/5+\eps}$. The existing generic benchmark is
$D^{16/15+\eps}$, leaving a $4/15$ power gap. A proof of the
proposed full moment would give
$\beta(2/5,2/5)=13/15<7/8$ by
Theorem~\ref{thm:finite-extraction}; no such proof is supplied here.
The parameter choice in Corollary~\ref{cor:packet-energy} proves
only its displayed controlled-sector estimate.
\end{remark}

\section{Signed gcd recombination at the original row scale}
\label{sfac:sec:gcd}\label{sec:gcd}

Set $x_u(n)=\mu_K(n)\lambda_u(n)W(\N n/D)$ and define
\begin{align}
 C(D,H)&=D^{-1}\sum_{u\ne0}\Phi(\N u/H)
       \sum_{\substack{n_1\ne n_2\\(n_1,n_2)=1}}
          x_u(n_1)\overline{x_u(n_2)},\notag\\
 B_d(D,H)&=D^{-1}\sum_{u\ne0}\Phi(\N u/H)
              \left|\sum_{d\mid n}x_u(n)\right|^2.
 \label{sfac:eq:gcd-definitions}
\end{align}

The sum over both ordered columns shows that \(C(D,H)\) also equals
the conjugated-orientation core \(\mathfrak C(D,H)\) used in
Section~\ref{sec:overlap}; swapping the two columns identifies them.

\begin{proposition}[Exact identity and quantitative truncation]
\label{sfac:prop:gcd}
For large $D$, when the unit column is off the profile,
\begin{equation}
 C(D,H)=\sum_d\mu_K(d)B_d(D,H).
 \label{sfac:eq:gcd-exact}
\end{equation}
For $T\ge1$ and $H\ge1$,
\begin{equation}
 C(D,H)-\sum_{\N d\le T}\mu_K(d)B_d(D,H)
 \ll_\varepsilon D^\varepsilon
        \left[H+\min\{HD/T,D^2/T^2\}\right].
 \label{sfac:eq:gcd-tail}
\end{equation}
The estimate uses the native primitive conductor and Poisson formulas
for ratios of coprime squarefree sextic columns.
\end{proposition}

\begin{proof}
Expanding the squares, the coefficient of a column pair is
$\sum_{d\mid(n_1,n_2)}\mu_K(d)=\mathbf1_{(n_1,n_2)=1}$.
Its only possible surviving diagonal is $n_1=n_2=1$, which is absent.
This proves \eqref{sfac:eq:gcd-exact}; all column and divisor sums are finite.

For distinct squarefree columns let $g=(n_1,n_2)$ and $G=\N g$.
Their ratio character, after cancellation of $g$, is nonprincipal and
primitive at the good conductor, of norm
$Q=\N(n_1n_2)/G^2$.  The original kernel retains the common-prime mask:
\[
 K_H(n_1,n_2)=\sum_{u\ne0}\Phi(\N u/H)
  \chi_{n_1/g}(u)\overline{\chi_{n_2/g}(u)}
  \mathbf1_{(u,g)=1}.
\]
Trivial counting and masked primitive completion give
$|K_H|\ll_\varepsilon D^\varepsilon\min\{H,D/G\}$.
For a dyadic gcd block $G\le\N g<2G$, counting $g$ and then its two
cofactors gives $O(D^2/G)$ pairs.  After division by $D$, the absolute
block bound is
$D^\varepsilon\min\{HD/G,D^2/G^2\}$.

Put $\sigma_T(g)=\sum_{d\mid g,\,\N d\le T}\mu_K(d)$.
For $1<\N g\le T$ this is zero, while for $\N g>T$ it is bounded by
$\tau(g)\ll_\varepsilon D^\varepsilon$.  Thus the off-diagonal
truncation error is bounded by the dyadic high-gcd tail, proving the
minimum term in \eqref{sfac:eq:gcd-tail}.  The truncated diagonal is
not generally zero.  Its absolute value is bounded by
\[
 D^{-1}\sum_n |\mu_K(n)\nu(n)W(\N n/D)|^2|\sigma_T(n)|
       \sum_{u\ne0}\Phi(\N u/H)\mathbf1_{(u,n)=1}
 \ll_\varepsilon HD^\varepsilon,
\]
by ideal and lattice counting.  This proves the remaining term.
\end{proof}

For squarefree $d$, extraction of $d\mid n$ gives
\[
 \sum_{d\mid n}x_u(n)=\mu_K(d)\lambda_u(d)
 \sum_{(m,dS)=1}\mu_K(m)\lambda_u(m)W(\N m/(D/\N d)).
\]
Both the column mask $(m,d)=1$ and the row mask $(u,d)=1$ survive.
If $\mathfrak M_d(X,H)$ denotes the corresponding inner moment
normalized by $1/X$, then
\begin{equation}
 B_d(D,H)=(\N d)^{-1}\mathfrak M_d(D/\N d,H).
 \label{sfac:eq:gcd-scale}
\end{equation}
The row scale is still $H$; it does not become $(D/\N d)^h$.

For $H=D^h$ and $a\ge0$, either $T\ge D^{1-a}$ or
$T\ge D^{1-(h+a)/2}$ makes \eqref{sfac:eq:gcd-tail} affordable at
$D^{h+a+\varepsilon}$.  At $a=0$ the smaller useful cutoff is
$D^{1-h/2}$.  This original-row cutoff differs from the
$D^{1/2}$ overlap cutoff in the earlier transformed kernel.  Moreover
$B_1=\mathfrak M_W(D,H)$: a triangle bound on
\eqref{sfac:eq:gcd-exact} retains the very moment to be proved.

There is a finite local check of the signed compensation.  For
$\Re s,\Re t>1$ and $a_p=\lambda_u(p)$,
\[
 \sum_{(m,n)=1}\frac{\mu_K(m)\mu_K(n)\lambda_u(m)
          \overline{\lambda_u(n)}}{(\N m)^s(\N n)^t}
 =\prod_{p\notin S}(1-a_p(\N p)^{-s}-\overline{a_p}(\N p)^{-t}).
\]
If the columns are all squarefree divisors of a product of $k$ primes
with $a_p=1$, and the profile is replaced by $1$ in this finite test,
the total including the unit column is $(-1)^k$, whereas its distinct
prime-only part is $k(k-1)$.  This identity illustrates compensation;
it makes no annular asymptotic assertion.

% Section fragment prepared for Edward Baker, 8 October 2026.
% GPT-6 (Codex), inherited configuration; exact serving variant and effort
% label are not exposed. Same-model internal derivation and audit.
% Uses master macros \N,\OO,\muK,\ind,\rad,\Mfrak,\eps.
% Imported arithmetic/Poisson identities and positive valid-range moment
% are explicitly separated from the new kernel/counting deductions.

\section{Complete row kernels and the controlled overlap sector}
\label{sec:complete-kernels}\label{sec:overlap}

This section records a second reduction of the short-family problem.
It is an alternative to the zero-integral factorization reduction.
The remaining signed form is not estimated here.
Keep the original fixed data over \(K=\mathbb Q(\sqrt{-3})\).
Suppress the fixed profile subscript on \(\Mfrak_W\) in this section.
Ideals are prime to \(S\), and their fixed primary generators are used
in residue symbols. A star denotes this restriction together with
squarefreeness. Let \(\operatorname{supp}W\subset[c,C]\subset(0,\infty)\).
Choose the fixed nonnegative radial Schwartz weight \(\Phi\), with
\(\Phi\ge1\) on \([0,1]\) and \(\widehat\Phi\) smooth and compactly
supported. Use the normalized radial Fourier transform
\[
 \widehat F(|w|^2)=\frac2{\sqrt3}\int_{\mathbb C}
      F(|z|^2)e(-zw)\,dx\,dy,\qquad
 e(z)=\exp(4\pi i\,\operatorname{Im}z/\sqrt3).
\]
The row lattice has covolume one in this measure. Fourier inversion
gives \(\widehat{\widehat F}=F\) for radial functions.

\subsection{The full diagonal and the original coprime core}

Write \(x_u(n)=\muK(n)\nu(n)\chi_n(u)W(\N(n)/D)\).
The original normalized smooth moment is
\(\Mfrak(D,H)=D^{-1}\sum_u\Phi(\N(u)/H)|\sum_nx_u(n)|^2\).
For large \(D\), the zero row contributes nothing.

\begin{proposition}[Full diagonal removal]\label{prop:ov-diagonal}
For \(H\ge1\),
\begin{align*}
 \Mfrak(D,H)&=\mathfrak D(D,H)+\mathfrak O(D,H),\\
 \mathfrak D(D,H)&=\frac1D\sum_n|\muK(n)\nu(n)W(\N(n)/D)|^2
       \sum_u\Phi(\N(u)/H)\ind_{(u,n)=1}.
\end{align*}
where \(\mathfrak O\) contains exactly the distinct column pairs, and
\(\mathfrak D(D,H)\ll H\). Its original coprime part is
\begin{equation}\label{eq:ov-original-core}
 \mathfrak C(D,H)=\frac1D\sum_u\Phi(\N(u)/H)
   \sum_{\substack{n_1\ne n_2\\(n_1,n_2)=1}}
               \overline{x_u(n_1)}x_u(n_2).
\end{equation}
\end{proposition}
\begin{proof}
Expand the square in the displayed conjugated orientation.
The zero extension gives \(|\chi_n(u)|^2=\ind_{(u,n)=1}\)
on the squarefree support. Lattice counting bounds the inner row
sum by \(O_\Phi(H)\), and ideal counting bounds the weighted column
sum by \(O_W(D)\). This proves the diagonal bound.
Equation~\eqref{eq:ov-original-core} is a subform of the remaining
distinct pairs. It is not a positive submoment.
\end{proof}

The diagonal above includes every Fourier frequency of each diagonal
column pair. Subtracting only the first Poisson zero frequency would
leave diagonal terms in the dual form.

\subsection{The finite transformed representation}

We specify the imported arithmetic packet rather than reproduce its
ray-group construction. There is a fixed finite collection of actual
coefficients \(a_\xi(n)\), with \(|a_\xi(n)|=1\) on good squarefree
ideals, and fixed constants \(c_\xi\). For distinct coprime columns,
put \(\psi=\chi_{n_1}\overline{\chi_{n_2}}\). The paired-Gauss identity is
\begin{equation}\label{eq:ov-paired-gauss}
 \sum_\xi c_\xi a_\xi(n_1)\overline{a_\xi(n_2)}
 =\muK(n_1)\muK(n_2)\overline{\nu(n_1)}\nu(n_2)
                                      \gamma(\overline\psi).
\end{equation}
For a primitive character of conductor \(\mathfrak m\), define
\[
 \gamma(\psi)=\N(\mathfrak m)^{-1/2}
      \sum_{x\bmod\mathfrak m}\psi(x)e(x/\mathfrak m),
\]
using the fixed generator in the quotient. Thus
\(|\gamma(\psi)|=1\) and
\(\gamma(\psi)\gamma(\overline\psi)=\psi(-1)\).
These are the finite identities of \cite[Lemma~4.1]{OA5}.
Their source-correct conventions, including the angular factor in
\(a_\xi\), are recorded in \cite{N3}; no arbitrary coefficient family
is substituted for them.

Put \(W_0(t)=t^{-1/2}\overline{W(t)}\). Primitive Poisson and the
finite changes of variables in \cite[Section~4.4]{OA5} give
\begin{equation}\label{eq:ov-transformed}
 \Mfrak-\mathfrak D=\sum_\xi c_\xi\mathfrak S^\circ_\xi,
\end{equation}
where
\begin{align}
 \mathfrak S^\circ_\xi
 ={}&\frac H{D^2}\!
 \sum_{\substack{b,f\\(b,f)=1}}^*\muK(f)\N(b)
 \sum_{\substack{m_1,m_2\\(m_1m_2,b)=1\\m_1\ne m_2}}^*
 a_\xi(m_1)\overline{a_\xi(m_2)}
 \chi_{m_1}(f)^4\overline{\chi_{m_2}(f)^4}
 \notag\\[-2mm]
 &\hspace{12mm}\cdot
 W_0(\N(bfm_1)/D)\overline{W_0(\N(bfm_2)/D)}
                              \mathcal K_f(m_1,m_2),\label{eq:ov-S}\\
 \mathcal K_f(m_1,m_2)
 ={}&\sum_{k\in\OO\setminus\{0\}}
 \chi_{m_1}(k)\overline{\chi_{m_2}(k)}
 \widehat\Phi\!\left(\frac{H\N(k)}
                  {\N(f)^2\N(m_1)\N(m_2)}\right).
 \label{eq:ov-K}
\end{align}
The \(f\)-characters retain the mask \((m_1m_2,f)=1\); the
\(b\)-mask is displayed. No coprimality between \(m_1,m_2\) is
imposed. The complete \(k\)-sum must be retained.
The original condition \(n_1\ne n_2\) becomes exactly \(m_1\ne m_2\).
These are finite Poisson/algebra identities, valid for \(H>0\);
they do not extend the source's deep moment theorem beyond its range.
For the variable bijection and the corrected complex-profile
orientation, see \cite{N2,N3}.

\begin{proposition}[Exact second-Poisson involution]
\label{prop:ov-involution}
In \eqref{eq:ov-S}, restrict to \(b=f=1\) and
\((m_1,m_2)=1\). After the original finite \(c_\xi\)-combination,
this restricted form equals \(\mathfrak C(D,H)\) in
\eqref{eq:ov-original-core}.
\end{proposition}
\begin{proof}
Set \(q=\N(m_1m_2)\) and
\(\psi=\chi_{m_1}\overline{\chi_{m_2}}\). The columns are distinct,
so \(\psi\) is nonprincipal and primitive of conductor \(m_1m_2\).
There is no shared-prime mask. Poisson at scale \(q/H\), followed
by radial Fourier inversion, gives exactly
\[
 \mathcal K_1(m_1,m_2)
 =\frac{\sqrt q}{H}\gamma(\psi)
               \sum_{u\ne0}\overline{\psi(u)}\Phi(\N(u)/H).
\]
The physical factors satisfy
\[
 \frac H{D^2}\frac{\sqrt q}{H}
 W_0(\N(m_1)/D)\overline{W_0(\N(m_2)/D)}
 =\frac1D\overline{W(\N(m_1)/D)}W(\N(m_2)/D).
\]
Now apply \eqref{eq:ov-paired-gauss}. The two Gauss factors multiply
to \(\psi(-1)\), which is removed from the complete radial row sum
by \(u\mapsto-u\). The resulting expression is precisely
\eqref{eq:ov-original-core}. Thus the second transform alone is an
involution, not a shorter-scale estimate.
\end{proof}

\subsection{Uniform primitive completion and overlap counting}

\begin{lemma}[Complete primitive kernel]\label{lem:ov-kernel}
For distinct good squarefree \(m_1,m_2\), put
\[
 g=(m_1,m_2),\quad G=\N(g),\quad
 \mathfrak c=m_1m_2/g^2,\quad Q=\N(\mathfrak c),\quad
 R=\frac{\N(f)^2\N(m_1)\N(m_2)}H.
\]
Then
\[
 |\mathcal K_f(m_1,m_2)|\ll_\Phi\tau(g)\sqrt Q.
\]
In the column ranges of \eqref{eq:ov-S}, this implies
\begin{equation}\label{eq:ov-min-kernel}
 |\mathcal K_f(m_1,m_2)|
 \ll_{\Phi,\eps}D^\eps\min\{R,\sqrt Q\}.
\end{equation}
If \(Z=\N(f)^2G/H\ge1\), then for every fixed \(A>1\),
\begin{equation}\label{eq:ov-rapid-kernel}
 |\mathcal K_f(m_1,m_2)|
 \ll_{\Phi,A}\tau(g)\frac R{\sqrt Q}Z^{-A}.
\end{equation}
\end{lemma}
\begin{proof}
At a prime present in only one column, the local sextic character
has exponent \(1\) or \(-1\), hence conductor exponent one.
Thus the ratio is a nonprincipal primitive character \(\psi\) of
conductor \(\mathfrak c\), with \((\mathfrak c,g)=1\), and
\[
 \chi_{m_1}(k)\overline{\chi_{m_2}(k)}
       =\psi(k)\ind_{(k,g)=1}.
\]
Apply the masked primitive Poisson formula
\cite[Lemma~4.2]{OA5} with weight \(\widehat\Phi\).
For a divisor \(d\mid g\), the prefactor has modulus
\(R/(\N(d)\sqrt Q)\), and the frequency scale is
\(t_d=R/(\N(d)Q)\). The zero frequency vanishes.
Lattice counting and Schwartz decay give
\[
 t\sum_{\ell\in\OO\setminus\{0\}}|\Phi(t\N(\ell))|
       \ll_\Phi1\qquad(t>0).
\]
Each divisor therefore costs \(O_\Phi(\sqrt Q)\).
This proves the uniform bound. Compact support of
\(\widehat\Phi\) gives the additional bound \(O_\Phi(R)\)
at every positive \(R\): a nonempty row sum forces \(R\)
above a fixed positive constant. Absorb \(\tau(g)\) into
\(D^\eps\) to obtain \eqref{eq:ov-min-kernel}.
Finally \(t_d=\N(f)^2G^2/(H\N(d))\ge Z\).
For \(t\ge1\), arbitrary Schwartz decay gives
\(\sum_{\ell\ne0}|\Phi(t\N(\ell))|\ll_{\Phi,A}t^{-A}\).
Sum the corresponding Poisson divisor terms to obtain
\eqref{eq:ov-rapid-kernel}.
\end{proof}

\begin{proposition}[Controlled overlap blocks]\label{prop:ov-block}
Let \(B,F,G\) be dyadic lower endpoints for
\(\N(b),\N(f),\N((m_1,m_2))\), respectively. The block of
\eqref{eq:ov-S} satisfies
\begin{equation}\label{eq:ov-block}
 |\mathfrak S^\circ_{\xi;B,F,G}|
 \ll_\eps D^\eps
       \min\left\{\frac{D^2}{B^2FG},
                    \frac{HD}{BF^2G^2}\right\}.
\end{equation}
The bound remains valid after a selector of modulus at most one
on \(b,f,m_1,m_2\) independent of \(k\). In particular, for
\(H=D^h\), all blocks with \(BF^2G^2\ge D^{1-a}\) contribute
\(O_\eps(D^{h+a+\eps})\).
\end{proposition}
\begin{proof}
The profile gives \(X=D/(BF)\), \(\N(m_j)\asymp X\),
\(R\asymp D^2/(HB^2)\), and \(\sqrt Q\asymp X/G\).
There are \(O(G)\) possible common ideals, and
\(O(X/G)\) choices for each quotient column, giving
\(O(X^2/G)\) pairs. Fixed support constants cover bounded
nonempty quotient ranges; impossible blocks are empty.
The counts of \(b,f\) and the exterior factor give
\[
 \frac{HB^2F}{D^2}\frac{X^2}{G}
       D^\eps\min\left\{\frac{D^2}{HB^2},\frac XG\right\},
\]
which simplifies to \eqref{eq:ov-block}. This estimates an
absolute majorant after each complete \(k\)-kernel has been assembled,
proving the selector statement. The dyadic block count is logarithmic,
and a smaller preliminary epsilon absorbs it.
\end{proof}

The cutoff is stated in dyadic lower endpoints. An actual norm product
\(\N(b)\N(f)^2\N(g)^2\ge32D^{1-a}\) necessarily belongs to a
controlled block. Boundary blocks may remain in the residual.
The estimate does not authorize a selector inside the \(k\)-sum.

\subsection{The imported valid range and the signed residual}

The only deep analytic input used in this sector removal is the
following valid-range source estimate \cite[Proposition~5.1]{OA5}.
For fixed \(\kappa>0,C_0\ge1\), \(R,X,F\ge1\),
\(\Sigma=XF\le D^{C_0}\), and \(R\le\Sigma D^{-\kappa}\),
\begin{equation}\label{eq:ov-imported-positive}
 \frac1{XF}\sum_{\N(f)\asymp F}^{*}
 \sum_{0<\N(k)\le R}
 \left|\sum_n^*a_\xi(n)\chi_n(k)\chi_n(f)^4V(\N(n)/X)\right|^2
 \ll_\eps \|V\|_{C^J}^2D^\eps\Sigma .
\end{equation}
Here \(J\) and the constant have the dependencies in the source.
The source's exclusion-removal comparison preserves \(\Sigma\)
and the row range. Smooth separation on fixed compact boxes and
Cauchy--Schwarz apply this estimate to the whole coupled quadratic
block of \eqref{eq:ov-S}; the resulting cost is \(HD^\eps\).
Indeed its outside factor and the two source norms give
\((HB^2/D^2)\Sigma^2=H\).
This application is made before imposing pair selectors.

Fix \(0<h<1\), \(a\ge0\), \(H=D^h\), and \(0<\eta<h/10\).
For complete blocks \(B\ge D^{1-h+\eta}\),
\[
 \frac R\Sigma\ll\frac{D}{HB}\ll D^{-\eta},
\]
so \eqref{eq:ov-imported-positive} applies with, for example,
\(\kappa=\eta/2\) after absorbing fixed constants.
Reintroducing and then subtracting the column diagonal costs
\(D/B\ll HD^{-\eta}\). Bounded nonempty scales not covered by
the source's \(R,X,F\ge1\) hypotheses are handled by direct counting.
Next, on the remaining blocks,
\eqref{eq:ov-rapid-kernel} removes
\(\N(f)^2\N(g)\ge HD^{2\eta}\) with error \(O(D^{-L})\)
for every requested \(L\). An explicit crude block cost is
\(D^{2+\eps-2\eta A}\), so choosing \(A\) first suffices.
Finally apply Proposition~\ref{prop:ov-block}.

Consequently the remaining signed form can be restricted to
\begin{equation}\label{eq:ov-residual}
 B<D^{1-h+\eta},\qquad
 \N(f)^2\N(g)<HD^{2\eta},\qquad
 BF^2G^2<D^{1-a},\qquad m_1\ne m_2,
\end{equation}
with all original coefficients, masks, profiles and complete rows.
Here the first and third restrictions concern dyadic lower endpoints,
removing any boundary ambiguity. If \(\mathfrak S_{\xi,\rm rem}\)
denotes this restriction, then
\[
 \Mfrak(D,D^h)=O_\eps(D^{h+a+\eps})
       +\sum_\xi c_\xi\mathfrak S_{\xi,\rm rem}+O(D^{-L}).
\]
An upper bound \(D^{h+a+\eps}\) for the real part of the retained
exterior combination therefore suffices for the original moment.
No such signed bound is proved here.
In particular \(B=F=G=1\) remains; the uniform block bound there
is only \(O(HD^{1+\eps})\).
The exact involution above explains why a second transform by itself
does not settle that core. The earlier positive long-row and
prime-selected counterchecks are documented in \cite{N2,N3,N7}.

\section{The imported generic sextic baseline}
\label{sec:generic-sextic-baseline}

We import de Faveri's sixth-order large sieve, rather than proving it.
For sixth-power-free good ideal indices and any coefficient vector,
its operator bound is
\begin{equation}\label{eq:imported-sextic-sieve}
 \sum_{\N v\le M}^{\text{$6$-free}}
 \left|\sum_{\N n\le N}^{\text{$6$-free}}b_n\chi_v^{\rm DF}(n)\right|^2
 \ll_\eps(MN)^\eps
 \left(M+N+M^{5/6}N^{1/3}+M^{1/3}N^{5/6}\right)\sum_n|b_n|^2.
\end{equation}
This is Theorem~1.1 of \emph{Optimal large sieve for fixed order characters},
\cite{DF}. Its finite character presentation and zero extension,
from Sections~2.3--2.5 of that source, are also part of the imported
input. In the Eisenstein PID one can choose its generators so that,
after finitely many unit and fixed bad-prime classes, a physical good
sixth-power-free row $v_0$ satisfies
$\chi_n(v_0)=\chi_n(\epsilon_v)\chi_v^{\rm DF}(n)$ with
$\epsilon_v\in\mu_6$. Splitting the six values of $\epsilon_v$ makes
the supplementary factor a fixed column twist. Both presentations
retain the same zeros at primes dividing $v$; no temporary decomposition
prime is inserted into the physical masks. On each class the column
coefficient vector is fixed while $v$ varies, so positivity allows
enlargement to the full row set in the imported operator bound.
The fixed-data construction and full ray-class comparison are
detailed in \cite{N6}.

\begin{corollary}[Physical operator on sixth-power-free columns]
\label{cor:physical-operator}
For a fixed coefficient vector supported on sixth-power-free good ideals
of norm at most \(X\), the imported sieve and stated presentation give
\begin{align}
 &\sum_{u\ne0}\Phi(\N u/H)
       \left|\sum_n b_n\chi_n(u)\right|^2\notag\\
 &\quad\ll_\eps(HX)^\eps
 [H+XH^{1/6}+H^{5/6}X^{1/3}+H^{1/3}X^{5/6}]
       \sum_n|b_n|^2. \label{eq:physical-operator}
\end{align}
\end{corollary}
\begin{proof}
For each fixed sixth-power factor \(r\), keep the column mask
\((n,r)=1\); its squared norm is at most \(\sum|b_n|^2\).
Apply the imported ideal-row sieve at \(H/(\N r)^6\).
The four terms sum with inverse powers \(6,0,5,2\) of \(\N r\).
The constant term counts \(O(H^{1/6})\) ideals and the other
three sums converge. Finite presentation splitting and Schwartz
row shells give the displayed bound.
\end{proof}

\begin{proposition}[Physical-row transfer of the generic sieve]
\label{prop:physical-generic-sieve}
Conditional on \eqref{eq:imported-sextic-sieve} and the stated
presentation, for $D,H\ge1$,
\begin{equation}\label{eq:physical-generic-sieve}
 \Mfrak_{\le}(D,H)\ll_\eps(DH)^\eps
 \left(H+DH^{1/6}+H^{5/6}D^{1/3}+H^{1/3}D^{5/6}\right).
\end{equation}
The same bound holds with a fixed Schwartz row weight. It also holds
for bounded squarefree-supported column coefficients in place of
$\muK\nu$.
\end{proposition}

\begin{proof}
Apply Corollary~\ref{cor:physical-operator} with \(X=CD\) and
\(b_n=\muK(n)\nu(n)W(\N n/D)\). Its squared norm is \(O_W(D)\)
by ideal counting. Divide by \(D\), and absorb the fixed support
constant. The same argument covers bounded squarefree-supported
coefficients and the sharp weight.
\end{proof}

For $1\le H\le D$ the right side is $O_\eps(D^{1+\eps}H^{1/6})$.
At $H=D^h$ it supplies only $a=1-5h/6$, whose extraction floor is one.
This power is necessary for the broad bounded coefficient class:
with trivial twist, a nonnegative nonzero annular profile, and prime-only
columns, every row $u=r^6$ with $\N r\le H^{1/6}$ has no prime mask on
the profile for large $D$, and has amplitude $\asymp D/\log D$ by the
prime ideal theorem. Ideal counting supplies $\asymp H^{1/6}$ such rows,
so their normalized energy is $\gg DH^{1/6}/(\log D)^2$.
This is a selected-coefficient obstruction, not a lower bound for the
unrestricted M\"obius family. A useful target \eqref{eq:improvement-budget}
requires more than a quarter power beyond that generic baseline and
must preserve additional arithmetic cancellation.

% Assembly references: notes 1, 5, 6; reviews
% SHORT_FAMILY_REFINEMENT_REVIEW_20261008.md and
% SHORT_FAMILY_CORE_AND_OVERLAP_REVIEW_20261008.md.

% Section fragment prepared for Edward Baker, 8 October 2026.
% Substantial GPT-6 (Codex) assistance; inherited configuration.
% Exact serving variant and configured effort are not exposed or inferred.
% Uses the manuscript's existing macros, theorem environments and labels.

\section{A nonsquarefree tail bound and a squarefree residual}
\label{sec:squarefree-tail}

The exact inverse creates nonsquarefree total products even though the
original M\"obius coefficient vanishes on them. Their complete contribution
can now be bounded, using the existing completion and physical operator.
This gives a squarefree residual at a different adaptive cutoff from
Section~\ref{sec:packets}; the signed moment of that residual remains open.
The detailed derivation is recorded in \cite{N17}.

Fix a zero-integral annular profile \(W\) and \(0<\theta<1\).
Write \(I_u^{(\theta)},T_u^{(\theta)}\) for the adaptive sums with
\(Y_u^{(\theta)}=D^{1-\theta}/L_u\). Define projection by the total ideal:
\begin{align}
 T_{u,\mathrm{nsq}}^{(\theta)}
 &=-\sum_{\substack{d,m\in\good\\
             \N d>Y_u^{(\theta)}\\dm\ \mathrm{nonsquarefree}}}
       c_z(d)\lambda_u(dm)W(\N(dm)/D),\notag\\
 \mathcal E_{\mathrm{nsq}}^{(\theta)}(D,H)
 &=\frac1D\sum_{u\ne0}\Phi(\N u/H)
       |T_{u,\mathrm{nsq}}^{(\theta)}|^2.\label{eq:nsq-projection}
\end{align}
The analogous small-product projection is denoted
\(I_{u,\mathrm{nsq}}^{(\theta)}\). Set
\[
 \mathcal B(H,X)=H+XH^{1/6}
                   +H^{5/6}X^{1/3}+H^{1/3}X^{5/6},
\]
the operator factor in Corollary~\ref{cor:physical-operator}.

\begin{theorem}[Complete nonsquarefree sector]\label{thm:nsq-tail}
For every fixed \(0<\theta<1\), every fixed smooth annular \(W\) with
\(\int W=0\), \(1\le H\le D\), and every \(\eps>0\),
\begin{equation}\label{eq:nsq-tail-bound}
 \mathcal E_{\mathrm{nsq}}^{(\theta)}(D,H)
 \ll_{\eps,\theta,W,\Phi}
 D^\eps\min\{HD^{1-\theta},
                    D^{-\theta/2}\mathcal B(H,D)\}.
\end{equation}
The second bound uses the imported physical operator on arbitrary
sixth-power-free columns. No new M\"obius or conductor-uniform
reciprocal estimate is assumed.
\end{theorem}

\begin{proof}
The exact factorization \eqref{sfac:eq:factor} and the vanishing of
\(\muK(n)\) on nonsquarefree ideals give the pointwise identity
\begin{equation}\label{eq:nsq-small-tail}
 T_{u,\mathrm{nsq}}^{(\theta)}
       =-I_{u,\mathrm{nsq}}^{(\theta)}.
\end{equation}
Use the finite divisor identity
\[
 \ind_{n\ \mathrm{nonsquarefree}}
       =-\sum_{\substack{k^2\mid n\\k\ne1}}\muK(k),
\]
and split the auxiliary ideals at \(R=D^\rho\), with
\(0<2\rho<\theta\). Let \(J_{u,\le R}\) and \(J_{u,>R}\) be the
resulting contributions to \(I_{u,\mathrm{nsq}}^{(\theta)}\).

For fixed \(d,k\), put \(j=k^2/(k^2,d)\). Ideal valuations give
\(k^2\mid dm\) if and only if \(j\mid m\), so complete multiplicativity,
including its zeros, gives
\[
 \sum_{m:k^2\mid dm}\lambda_u(m)W(\N(dm)/D)
 =\lambda_u(j)S_u\!\left(\frac D{\N d\,\N j}\right).
\]
For \(\N d\le Y_u^{(\theta)}\) and \(\N k\le R\),
\(\N j\le R^2\) and \(L_u\N d\,\N j/D\le D^{-\theta}R^2\).
There are \(O(R)\) auxiliary ideals, and
\(\sum_{\N d\le Y}|c_z(d)|\ll Y\log(2Y)\).
Lemmas~\ref{sfac:lem:completion} and~\ref{sfac:lem:mass} therefore imply,
for every completion order \(M\ge0\) and \(\delta>0\),
\begin{equation}\label{eq:nsq-small-square}
 \frac1D\sum_{u\ne0}\Phi(\N u/H)|J_{u,\le R}|^2
 \ll_{M,\delta,W,\Phi}
 D^{1-2\theta}R^2\log^2(2D)H^\delta
                         (D^{-\theta}R^2)^{2M}.
\end{equation}
The small sum is empty when \(Y_u^{(\theta)}<1\).
Since \(\theta-2\rho>0\) and \(H\le D\), this energy is
\(O_N(D^{-N})\) for every requested \(N\), choosing \(M\) afterwards.

For the remaining part, define the fixed coefficients
\[
 w_R(n)=-\sum_{\substack{k^2\mid n\\\N k>R}}\muK(k),
 \qquad
 v_Y(n)=-\sum_{\substack{d\mid n\\\N d\le Y}}c_z(d).
\]
Then
\[
 J_{u,>R}=\sum_n w_R(n)v_{Y_u^{(\theta)}}(n)
                         \lambda_u(n)W(\N n/D).
\]
With \(\tau_3=\mathbf1*\mathbf1*\mathbf1\), uniformly in \(Y\),
\[
 |v_Y(n)|\le\tau_3(n),\qquad
 |w_R(n)|\le\sum_{\substack{k^2\mid n\\\N k>R}}1.
\]
The divisor bound and ideal counting yield
\[
 |J_{u,>R}|
 \ll_{\eps,W}D^\eps
       \sum_{\N k>R}\#\{m:\N(k^2m)\le CD\}
 \ll_{\eps,W}D^{1+\eps}/R,
\]
because \(\sum_{\N k>R}(\N k)^{-2}\ll R^{-1}\).
The full row weight has mass \(O_\Phi(H)\), hence
\begin{equation}\label{eq:nsq-elementary-large}
 \frac1D\sum_{u\ne0}\Phi(\N u/H)|J_{u,>R}|^2
                  \ll_{\eps,W,\Phi}HD^{1+\eps}/R^2.
\end{equation}

For a sharper bound, order the good ideals \(d\) of norm at most \(CD\)
by norm, fixing ties. The condition \(\N d\le Y_u^{(\theta)}\) is a
prefix. A fixed binary partition expresses each prefix as
\(O(\log D)\) intervals, with \(O(\log D)\) partition levels.
Thus a maximal-prefix bound follows by Cauchy and the sum of squares
over the fixed intervals; no row-dependent vector is passed to an
operator.

Uniquely write \(n=\ell^6n_0\), with \(n_0\) sixth-power-free, and put
\(X_\ell=CD/(\N\ell)^6\). For a fixed partition interval \(\mathcal J\),
the column vector, up to its common sign, is
\[
 b_{\mathcal J,\ell}(n_0)=
 w_R(\ell^6n_0)\nu(\ell^6n_0)W(\N(\ell^6n_0)/D)
       \sum_{\substack{d\mid\ell^6n_0\\d\in\mathcal J}}c_z(d).
\]
The physical identity
\(\chi_{\ell^6n_0}(u)=\ind_{(u,\ell)=1}\chi_{n_0}(u)\)
retains the common row mask. For each \(\ell\) it may be removed only
by enlarging a nonnegative sum of squares.

If \(w_R(\ell^6n_0)\ne0\), some squarefree \(k\), \(\N k>R\),
has \(k^2\mid\ell^6n_0\). Set \(g=(k,\ell)\), \(k_2=k/g\).
Squarefreeness of \(k\) gives
\[
 k_2^2\mid n_0,\qquad \N k_2>R/\N\ell.
\]
Summing ideal counts over these square divisors, or using the
unrestricted count when \(R/\N\ell<1\), shows that the number of
possible columns is
\[
 \ll X_\ell\min\{1,\N\ell/R\}.
\]
On each partition level the intervals are disjoint in \(d\), so
\[
 \sum_{\mathcal J}
 \left|\sum_{\substack{d\mid n\\d\in\mathcal J}}c_z(d)\right|^2
 \le\left(\sum_{d\mid n}|c_z(d)|\right)^2\le\tau_3(n)^2.
\]
Also \(|w_R(n)|\le\tau(n)\). Consequently, on each level,
\[
 \sum_{\mathcal J}\norm{b_{\mathcal J,\ell}}_2^2
 \ll_{\eps,W}D^\eps X_\ell\min\{1,\N\ell/R\}.
\]
Corollary~\ref{cor:physical-operator}, applied to these fixed
sixth-power-free vectors, bounds the \(\ell\)-channel with any prefix
by this quantity times \(D^\eps\mathcal B(H,X_\ell)\).
All partition and maximal-prefix logarithms are absorbed by
reallocating epsilon.

Finally weighted Cauchy in \(\ell\) uses
\(\sum_\ell(\N\ell)^{-2}<\infty\). After division by \(D\),
\begin{align}
 \frac1D\sum_{u\ne0}\Phi(\N u/H)|J_{u,>R}|^2
 &\ll_{\eps,W,\Phi}D^\eps
   \sum_\ell(\N\ell)^{-4}\min\{1,\N\ell/R\}
                       \mathcal B(H,X_\ell)\notag\\
 &\ll_{\eps,W,\Phi}\frac{D^\eps}{R}\mathcal B(H,D)
                         \sum_\ell(\N\ell)^{-3}
 \ll_{\eps,W,\Phi}D^\eps\mathcal B(H,D)/R.
 \label{eq:nsq-sparse-large}
\end{align}
Only \(X_\ell\ge1\) can contribute. The fixed annular constant \(C\)
is absorbed when comparing the operator factors.

Choose \(\rho=\theta/2-\kappa\), where the fixed positive \(\kappa\)
and the preliminary divisor and operator epsilons are sufficiently
small in terms of the requested \(\eps,\theta\).
Equations~\eqref{eq:nsq-small-square}--\eqref{eq:nsq-sparse-large}
then give \eqref{eq:nsq-tail-bound}. Both large-square bounds estimate
the same vector, so taking their minimum is legitimate.
\end{proof}

\begin{corollary}[An affordable squarefree residual]\label{cor:sf-residual}
Let \(H=D^{2/5}\), \(\theta=11/20\), and \(\int W=0\). Then
\begin{equation}\label{eq:nsq-affordable}
 \mathcal E_{\mathrm{nsq}}^{(11/20)}(D,D^{2/5})
       \ll_{\eps,W,\Phi}D^{19/24+\eps},
 \qquad \frac45-\frac{19}{24}=\frac1{120}.
\end{equation}
If \(T_{u,\mathrm{sf}}^{(\theta)}\) retains only squarefree total
ideals, the full moment target
\(\mathfrak M_W(D,D^{2/5})\ll_\eps D^{4/5+\eps}\)
is equivalent to
\begin{equation}\label{eq:sf-residual-target}
 \frac1D\sum_{u\ne0}\Phi(\N u/D^{2/5})
       |T_{u,\mathrm{sf}}^{(11/20)}|^2
                       \ll_\eps D^{4/5+\eps}.
\end{equation}
The estimate \eqref{eq:sf-residual-target} remains open.
\end{corollary}

\begin{proof}
At \(H=D^{2/5}\), \(\mathcal B(H,D)\ll D^{16/15}\), and
\(16/15-11/40=19/24\). Theorem~\ref{sfac:thm:adaptive} equates the
full zero-integral moment with its full adaptive tail moment.
Equation~\eqref{eq:nsq-affordable} then permits removing or restoring
the nonsquarefree response by the weighted triangle inequality.
\end{proof}

In the remaining factorization, squarefreeness of the total ideal
\(n=abm\) forces \(a,b,m\) to be squarefree and pairwise coprime.
These restrictions are now justified with a proved error budget;
they were not permissible in the original exact convolution.
The restricted \(m\)-sum is no longer a free ideal sum, so its
completion requires a separate argument.

Since \(Y_u^{(11/20)}\le D^{9/20}<z\) for sufficiently large \(D\),
\eqref{eq:packet-inverse-coefficient} also gives the exact residual form
\[
 T_{u,\mathrm{sf}}^{(11/20)}
 =\sum_{n\ \mathrm{squarefree}}
 \left[\muK(n)+
   \sum_{\substack{d\mid n\\\N d\le D^{9/20}/L_u}}
                    (-2)^{\omega(d)}\right]
              \lambda_u(n)W(\N n/D).
\]
The inner sum is empty when \(Y_u^{(11/20)}<1\). This identity
does not estimate the remaining signed response.

The cutoff differs from \(\theta=1/40\) in
Corollary~\ref{cor:packet-energy}. For zero-integral profiles, the
full vectors at any two fixed positive buffers differ by an
arbitrarily small weighted response, by
Theorem~\ref{sfac:thm:adaptive}; their product projections need not.
The controlled sectors from the two representations cannot therefore
be combined as projections of one unchanged tail.

\subsection{Complete cofactors and odd rough cores}
\label{sec:rough-parity}

The following finite decomposition, developed in~\cite{N21}, uses the
same buffer \(\theta=11/20\) as Corollary~\ref{cor:sf-residual}.
It does not transfer a projected estimate from another buffer.
For a squarefree total ideal \(n\) and \(1\le Y<z\), put
\[
 b_Y(n)=\prod_{\substack{p\mid n\\\N p\le Y}}p,
 \qquad r_Y(n)=n/b_Y(n),\qquad
 O_Y(b)=\sum_{\substack{d\mid b\\\N d>Y}}(-2)^{\omega(d)}.
\]
Both factors are squarefree and coprime. Prime ideals of equal norm
enter the first factor simultaneously.

\begin{proposition}[Complete cofactor parity]\label{prop:rough-parity}
For squarefree \(n\) on the profile, the actual tail coefficient satisfies
\begin{equation}\label{eq:rough-parity}
 t_Y(n):=-\sum_{\substack{d\mid n\\\N d>Y}}c_z(d)
  =\muK(b_Y)\{1+\muK(r_Y)\}-O_Y(b_Y).
\end{equation}
If \(\N b_Y\le Y\), this coefficient is zero when
\(\omega(r_Y)\) is odd and equals \(2\muK(b_Y)\) when it is even.
The unit rough core has even parity. If \(Y<1\), instead
\(t_Y(n)=\muK(n)\).
\end{proposition}
\begin{proof}
Every divisor with norm at most \(Y\) divides \(b_Y\). Also
\[
 \sum_{d\mid b_Y}(-2)^{\omega(d)}
       =\prod_{p\mid b_Y}(1-2)=\muK(b_Y).
\]
Insert this identity, retaining its omitted divisors \(O_Y\), in the
squarefree coefficient formula above. On \(\N b_Y\le Y\) the omitted
sum is empty. For \(Y<1\) the original prefix is empty.
All tuples at a total ideal carry the same \(\lambda_u(n)W(\N n/D)\),
so the cancellation retains every complex phase and deletion zero.
\end{proof}

This gives a new nonempty signed sector at growing cofactor width.
Fix \(\kappa>0\), \(0<\delta<11/40\), and let \(\mathcal Q(D)\)
be the set of squarefree total ideals on the profile represented by
\[
 n=bq_1q_2q_3,\qquad \N b\le D^{1/40},\qquad
 \N q_j\ge\kappa D^{13/40},
\]
where the \(q_j\) are distinct good primes and \((b,q_1q_2q_3)=1\).
Each total ideal is counted once. Define
\[
 C_u=\ind_{L_u\ge D^{1/8+\delta}}
       \sum_{n\in\mathcal Q(D)}t_{Y_u}(n)\lambda_u(n)W(\N n/D),
 \qquad Y_u=D^{9/20}/L_u.
\]

\begin{theorem}[A growing odd-core sector at large conductor]
\label{thm:rough-sector}
For every \(N>0\), under the existing conductor comparison and complete
Schwartz weight,
\begin{equation}\label{eq:rough-sector-energy}
 \frac1D\sum_{u\ne0}\Phi(\N u/D^{2/5})|C_u|^2
       =O_{N,\delta,\kappa,W,\Phi}(D^{-N}).
\end{equation}
\end{theorem}
\begin{proof}
On \(\N u\le D^{33/80}\), the fixed-constant bound \(L_u\ll\N u\)
gives
\[
 Y_u\gg D^{3/80}>D^{1/40}\ge\N b.
\]
For rows selected in \(C_u\), also
\[
 Y_u\le D^{13/40-\delta}<\kappa D^{13/40}\le\N q_j.
\]
These strict comparisons hold for sufficiently large \(D\), absorbing
all fixed constants. Thus \(b=b_{Y_u}(n)\) is complete below the
prefix cutoff and the rough core has three primes.
Proposition~\ref{prop:rough-parity} makes every selected coefficient
zero on this row range. Uniformly on the exterior rows, divisor bounds
and ideal counting give \(|C_u|\ll_{\eps,W}D^{1+\eps}\). Schwartz
decay gives, for every fixed \(A>0\),
\[
 \sum_{\N u>D^{33/80}}\Phi(\N u/D^{2/5})
     \ll_{A,\Phi}D^{2/5-A/80}.
\]
The exterior normalized energy is
\(O(D^{7/5+2\eps-A/80})\). Choose \(A\) after \(N\).
\end{proof}

The sector has growing cofactor width: fixed disjoint prime windows
near \(\alpha D^{1/40}\) and \(\beta_jD^{13/40}\), with
\(0<\alpha<1\), \(\beta_j>\kappa\), and their product inside the
annulus, contain \(\asymp D/\log^4D\) distinct total ideals by the
already imported fixed-field prime ideal theorem. This counting input
is not needed for the sector estimate.
This is cancellation of actual tail terms. For example, when \(b=1\)
and \(q_j\asymp D^{1/3}\), the singleton coefficients of \(c_z\)
are \(-2\), the three pair coefficients are \(+2\), and the triple
coefficient is zero. On \(1\le Y_u<\min_j\N q_j\), the singleton
and pair tail contributions are \(+6\) and \(-6\).
By contrast, on coherent inner rows \(u=r^6\),
\(L_u\ll\N r\le D^{1/15}\), so \(Y_u\gg D^{23/60}\), and
balanced triples retain coefficient \(-6\). Balanced semiprimes retain
coefficient \(+2\) on every row with \(Y_u\ge1\).
Rows with \(Y_u<1\) require \(L_u>D^{9/20}\), hence
\(\N u\gg D^{9/20}\), and have negligible full weighted energy at
\(H=D^{2/5}\). No low-conductor sector is silently removed.

For clarity about the remaining budget, set
\[
 g_Y(n)=2\muK(b_Y)\ind_{\N b_Y\le Y}
                  \ind_{\omega(r_Y)\ \mathrm{even}},\qquad
 h_Y(n)=t_Y(n)-g_Y(n)
\]
when \(Y\ge1\), and set \(g_Y=0,h_Y=\muK\) when \(Y<1\).
Let \(G_u,R_u\) be their squarefree responses. Then exactly
\(T_{u,\mathrm{sf}}^{(11/20)}=G_u+R_u\): the first term retains
saturated even rough cores and the second retains unsaturated cofactors.
The required moment keeps \(G_u+R_u\) inside one absolute square.
Each canonical response has only the generic bound
\begin{equation}\label{eq:rough-retained-generic}
 \frac1D\sum_{u\ne0}\Phi(\N u/H)(|G_u|^2+|R_u|^2)
       \ll_{\eps,W,\Phi}D^\eps B(D,H).
\end{equation}
Indeed, extend \(t_y\) on the auxiliary common threshold list by the
original truncated tail at every \(y\); its low-prefix formula is only
used at the actual \(Y_u<z\). For each fixed column, its total jump variation is at
most \(3^{\omega(n)}\), and \(g_y\) has \(O(1+\omega(n))\)
variation, with changes only at divisor norms. The same bound up to
a divisor loss holds for \(h_y\). Express their common integer-norm
threshold prefixes in fixed binary intervals. At one partition level
the sum of squared interval coefficients for a column is bounded by
the square of its total variation, hence by \(O_\eps(D^\eps)\).
Apply Corollary~\ref{cor:physical-operator} to each fixed interval
vector, sum over the \(O(D)\) squarefree columns and the logarithmic
levels, and divide by \(D\). The constant initial vector has the same
bound. This proves \eqref{eq:rough-retained-generic} without inserting
a row selector into a completed kernel.
At \(H=D^{2/5}\) it still costs \(D^{16/15+\eps}\).
The new zero sector therefore gives no improved full-moment exponent.

\section{Quantitative cutoffs for coherent cores}
\label{sfac:subsec:cores}

The decomposition $u=\eta vr^6$ also permits a precise separation of
fixed-core information from growing-conductor assumptions.  Use the
finite unit and bad-prime splitting, retain the complete row weight,
and put
\[
 \tau_{\eta,v}(n)=\nu(n)\chi_n(\eta v),\qquad
 C_{\eta,v}(X)=\sum_{n\in\good}
      \mu_K(n)\tau_{\eta,v}(n)W(\N n/X).
\]
Euler-factor completion gives the exact masked identity
\begin{equation}
 A_{\eta vr^6,W}(D)=
   \sum_{\substack{d\in\good\\
                  \operatorname{rad}(d)\mid\operatorname{rad}(r)}}
       \tau_{\eta,v}(d)C_{\eta,v}(D/\N d).
 \label{sfac:eq:core-completion}
\end{equation}
Every good prime power supported on $\operatorname{rad}(r)$ is allowed
in this sum; primes in $S$ have no effect on the completion.  Zeros at primes dividing $vS$ are retained; the annular
profile makes the sum finite.

\begin{proposition}[Conditional low-core estimate]
\label{sfac:prop:core-cutoff}
Fix $\beta>0$.  Assume, explicitly, that for each core under
consideration and every $\varepsilon>0$,
\begin{equation}
 |C_{\eta,v}(X)|\ll
 C_{\eta,v,\varepsilon}X^{\beta+\varepsilon}
 \label{sfac:eq:core-hypothesis}
\end{equation}
at all nonempty scales, including bounded scales.  For $1\le V\le H$,
the moment restricted to cores $\N v\le V$ satisfies
\begin{equation}
 \mathfrak M_{\le V}(D,H)\ll_\varepsilon
 D^{2\beta-1+\varepsilon}H^{1/6}
       \sum_\eta\sum_{\N v\le V}
           C_{\eta,v,\varepsilon}^{\,2}(\N v)^{-1/6}.
 \label{sfac:eq:low-core}
\end{equation}
In particular, an additional conductor-uniform hypothesis
$C_{\eta,v,\varepsilon}\ll_\varepsilon(\N v)^\kappa$, $\kappa\ge0$,
would give
\begin{equation}
 \mathfrak M_{\le V}(D,H)\ll_\varepsilon
   D^{2\beta-1+\varepsilon}H^{1/6}V^{2\kappa+5/6}.
 \label{sfac:eq:low-core-uniform}
\end{equation}
Fixed-character estimates alone do not provide this uniformity.
\end{proposition}

\begin{proof}
Use a smaller preliminary epsilon in \eqref{sfac:eq:core-hypothesis}.
Equation \eqref{sfac:eq:core-completion} bounds each row amplitude by
$C_{\eta,v,\varepsilon}D^{\beta+\varepsilon}F_\beta(r)$, where
\[
 F_\beta(r)=\prod_{p\mid r}(1-(\N p)^{-\beta})^{-1}.
\]
For this auxiliary row-ideal average, convolution is on all integral
ideals, including primes in $S$, rather than on the coefficient monoid.
The multiplicative function $F_\beta^2=\mathbf1*g_\beta$ has
nonnegative squarefree-supported coefficients with
$g_\beta(p)=(1-(\N p)^{-\beta})^{-2}-1$.  Its Euler product
\[
 \sum_d\frac{g_\beta(d)}{\N d}
   =\prod_p\left(1+
    \frac{(1-(\N p)^{-\beta})^{-2}-1}{\N p}\right)
\]
converges because $\beta>0$.  Ideal counting consequently gives
$\sum_{\N r\le R}F_\beta(r)^2\ll_\beta R$ for $R\ge1$.
Dyadic shells against $\Phi(\N v(\N r)^6/H)$ have total
$O((H/\N v)^{1/6})$ when $\N v\le H$.  Summing the squared amplitudes
and dividing by $D$ proves \eqref{sfac:eq:low-core}.
Under the additional uniformity, ideal counting gives
$\sum_{\N v\le V}(\N v)^{2\kappa-1/6}\ll V^{2\kappa+5/6}$.
\end{proof}

For comparison, the imported ideal-row sieve
\eqref{eq:imported-sextic-sieve}, with its stated native-family transfer,
has operator factor
\[
 \Theta_6(U,X)\ll_\varepsilon(UX)^\varepsilon
   [U+X+U^{5/6}X^{1/3}+U^{1/3}X^{5/6}].
\]
Keeping each fixed $r$-mask, its complementary-core consequence is
\begin{equation}
 \mathfrak M_{>V}(D,H)\ll_\varepsilon(DH)^\varepsilon
 \left[H+D(H/V)^{1/6}
          +H^{5/6}D^{1/3}+H^{1/3}D^{5/6}\right].
 \label{sfac:eq:high-core-generic}
\end{equation}
To see the multiplicity budgets, in a bounded row shell fix $r$ and
apply the sieve at core length $H/(\N r)^6$.  Only
$\N r\ll(H/V)^{1/6}$ can pay its $D$ term.  The $H$, $H^{5/6}$,
and $H^{1/3}$ terms sum respectively with inverse powers
$(\N r)^{-6}$, $(\N r)^{-5}$, and $(\N r)^{-2}$, all convergent.
Schwartz row shells recover the full weight.  This deduction retains
the imported status of the sieve and its native-family transfer.

For example, even granting $\beta=7/8$, a fixed finite core collection
has the supplied upper-bound exponent $3/4+h/6$; it does not improve
that benchmark scalar exponent.  Moreover the last term of
\eqref{sfac:eq:high-core-generic} has exponent $5/6+h/3$, larger
than every useful budget $h+a<3/4+h/6$ by more than
$1/12+h/6$.  This is an insufficiency of the supplied upper bound,
not a lower bound for the actual Möbius energy.



\section{Selected pieces that cannot meet the useful budget}
\label{sec:obstructions}

The following obstructions concern actual selected components.
They rule out giving the target bound separately to every such
component. They do not lower-bound the full M\"obius moment.

\subsection{The original prime-selected coprime core}

\begin{proposition}[Prime-column obstruction]\label{prop:prime-core}
Let \(\nu=1\), let \(W\) be nonnegative and nonzero, and let
\(H=D^h\), \(0<h<1\).
The coprime off-diagonal form restricted to distinct good prime
columns satisfies
\begin{equation}\label{eq:prime-core-lower}
 \mathfrak C_{\rm prime}(D,H)
 \gg \frac{DH^{1/6}}{(\log D)^2}-O\!\left(\frac H{\log D}\right)
 \gg \frac{DH^{1/6}}{(\log D)^2}.
\end{equation}
\end{proposition}
\begin{proof}
Put \(P_u(D)=\sum_{p\notin S}\chi_p(u)W(\N p/D)\).
The distinct-prime form is
\[
 \frac1D\sum_{u\ne0}\Phi(\N u/H)|P_u(D)|^2
 -\frac1D\sum_p|W(\N p/D)|^2
          \sum_{u\ne0}\Phi(\N u/H)\ind_{(u,p)=1}.
\]
Indeed \(\muK(p)=-1\), and distinct prime ideals are coprime.
The diagonal is \(O(H/\log D)\) by lattice and prime ideal counting.
There are \(\asymp H^{1/6}\) distinct coherent rows \(u=r^6\) in
the inner ball. Every prime on the profile has norm \(\asymp D\);
none divides \(r\), whose norm is at most \(H^{1/6}\).
Their \(P_u(D)\) are consequently all
\(\sum_pW(\N p/D)\asymp D/\log D\).
Positivity of the full square and the inner-ball weight prove the
first inequality. Its first term dominates for \(0<h<1\).
\end{proof}

The exact involution in Proposition~\ref{prop:ov-involution}
places this obstruction in the corresponding transformed coprime
core with its actual finite coefficient combination.
It requires \(a+5h/6\ge1\) for a separate bound
\(D^{h+a+\eps}\), while a useful target requires
\(a+5h/6<3/4\). The unrestricted core can compensate the prime
columns by other column types.

\subsection{A prime-by-prime component inside the adaptive tail}

Zero integral removes the principal frequency, but does not make
every remaining factor class small.
We now use the quantitative fixed-field prime ideal theorem
\begin{equation}\label{eq:quantitative-pnt}
 \pi_K(x)=\operatorname{Li}(x)
       +O_K(xe^{-c_K\sqrt{\log x}})
 \quad(x\text{ sufficiently large}).
\end{equation}
It follows from the fixed-field specialization of
\cite[Corollary~1.4]{DKN}, subtracting prime powers and using partial
summation. A possible fixed exceptional-zero term is absorbed after
increasing the threshold. The proof below needs this quantitative
remainder, rather than an unquantified asymptotic.

\begin{theorem}[Actual prime-by-prime tail obstruction]
\label{thm:prime-pair-tail}
Let \(\nu=1\), \(V\ge0\) be a nonzero smooth function supported
inside \((c,C)\), and \(W=(1+t\partial_t)V\).
For fixed \(0<\theta<1\), \(H=D^h\), \(0<h<1\), consider the
adaptive tail of Theorem~\ref{sfac:thm:adaptive}.
Select its ordered factor tuples
\[
 a=p,\quad b=q,\quad p\ne q,\quad
 \N p,\N q\le z=\sqrt{CD},\qquad m=1.
\]
Their actual amplitude is
\[
 U_u(D)=-\sum_{\substack{p\ne q\\\N p,\N q\le z}}
       \chi_p(u)\chi_q(u)W(\N p\N q/D),
\]
with good primes throughout. Every nonzero summand is in the adaptive
tail for sufficiently large \(D\), and
\begin{equation}\label{eq:prime-pair-tail-energy}
 \frac1D\sum_{u\ne0}\Phi(\N u/H)|U_u(D)|^2
       \gg_{V,\Phi}\frac{DH^{1/6}}{(\log D)^4}.
\end{equation}
\end{theorem}
\begin{proof}
The original coefficient on these tuples is
\(-\muK(p)\muK(q)=-1\).
Profile support and the two upper bounds force both prime norms
into a fixed annulus of size \(D^{1/2}\).
Apply \eqref{eq:quantitative-pnt} in each prime variable.
After \(\N p=D^{1/2}x\), \(\N q=D^{1/2}y\), the product of the
two prime densities is \(4/(\log D)^2+O((\log D)^{-3})\),
uniformly on this fixed compact region. Partial summation, including
the upper boundaries, therefore gives
\begin{align}
 \sum_{\N p,\N q\le z}W(\N p\N q/D)
 &=\frac{4D}{(\log D)^2}
     \int_0^{\sqrt C}\int_0^{\sqrt C}W(xy)\,dx\,dy
          +O_V(D/(\log D)^3)\notag\\
 &=\frac{4D}{(\log D)^2}
     \int_c^CW(t)\log(C/t)\,dt+O_V(D/(\log D)^3).
 \label{eq:prime-pair-asymptotic}
\end{align}
For the second identity, fix \(x\), put \(t=xy\), and integrate
\(dx/x\) between \(t/\sqrt C\) and \(\sqrt C\).
Since \(W=(tV)'\), integration by parts gives
\[
 \int_c^CW(t)\log(C/t)\,dt=\int_c^CV(t)\,dt>0.
\]
The omitted diagonal \(p=q\) is \(O_V(\sqrt D/\log D)\),
smaller than the error above.

For every coherent row \(u=r^6\), \(\N u\le H\), the prime
norms \(\gg\sqrt D\) exceed \(\N r\le D^{h/6}\);
both masks are one. Moreover
\(\N p\N q\ge cD>D^{1-\theta}/L_u=Y_u\), since \(L_u\ge1\).
Thus every such row has amplitude
\[
 U_{r^6}(D)
 =-\frac{4D}{(\log D)^2}\int V(t)\,dt
       +O_V(D/(\log D)^3).
\]
Count the \(\asymp H^{1/6}\) distinct coherent rows, square,
use \(\Phi\ge1\) on the inner ball, and divide by \(D\).
\end{proof}

The power in \eqref{eq:prime-pair-tail-energy} exceeds every useful
budget by more than \(1/4\), up to arbitrary small epsilon losses.
This is a lower bound for the selected amplitude only.
It neither lower-bounds the full \(m=1\) block nor permits removal
of the cross terms with composite factors and \(m>1\).
For the original prime columns on small-modulus rows, the large
primes occur in the removed \(d=1\) sector. Their zero-integral
cancellation is retained by the adaptive reduction. They must not
be used as a lower bound for its new tail; the theorem above gives
a separate tail-specific obstruction.

\subsection{Complete small-cofactor pairing retains the coherent power}

The preceding obstruction selects factor tuples. The next one, developed
in~\cite{N24}, first pairs actual total products with every divisor of
a growing small cofactor. Its signs include both semiprimes and triples.

\begin{theorem}[Complete small-cofactor obstruction]
\label{thm:complete-cofactor-obstruction}
Let \(\nu=1\), \(H=D^{2/5}\), \(z=\sqrt{CD}\), and
\(Y_u=D^{9/20}/L_u\). Fix a nonzero complex
\(W\in C_c^\infty((c,C))\) with \(\int W=0\), and fix
\(0<\beta<1/20\). For any squarefree good ideal \(M=M(D)\)
with \(\N M\le D^\beta\), define
\[
 \mathcal P_M=\{bqr:b\mid M,\ q\ne r\text{ good primes},\
                 \N q,\N r\le z,\ (qr,M)=1\},
 \qquad
 B_{u,M}=\sum_{n\in\mathcal P_M}
                  t_{Y_u}(n)\lambda_u(n)W(\N n/D),
\]
where each total ideal is counted once and \(t_Y\) is its actual
squarefree tail coefficient from \eqref{eq:rough-parity}.
Put \(P_-(M)=\prod_{p\mid M}(1-(\N p)^{-1})\). There is a
fixed integer \(k_0=k_0(W)\ge1\) such that, uniformly in \(M\),
\begin{equation}\label{eq:complete-cofactor-energy}
 \frac1D\sum_{u\ne0}\Phi(\N u/H)|B_{u,M}|^2
 \gg_{\beta,W,\Phi}
 \frac{D H^{1/6}P_-(M)^3}{(\log D)^{2k_0+2}}
 \gg_{\beta,W,\Phi}
 \frac{D^{16/15}}
      {(\log D)^{2k_0+2}(\log\log D)^3}.
\end{equation}
This lower bound concerns the selected signed response, not the full tail.
\end{theorem}

\begin{proof}
For \(u=v^6\), \(0<\N v\le D^{1/15}\), and \((v,M)=1\),
the radical comparison gives \(Y_{v^6}\gg D^{23/60}>\N M\).
Every contributing prime satisfies
\[
 \N q,\N r\ge cD/(\N b\,z)
       \gg D^{1/2-\beta}>D^{9/20}\ge Y_{v^6}.
\]
Thus all physical masks are one, the full small-prime part is \(b\),
and Proposition~\ref{prop:rough-parity} gives \(t_{Y_{v^6}}(bqr)=2\muK(b)\).
The ordered prime pairs absorb this factor two. Hence
\begin{equation}\label{eq:complete-cofactor-coherent}
 B_{v^6,M}=\sum_{b\mid M}\muK(b)
       \sum_{\substack{q\ne r\\\N q,\N r\le z}}
                     W(\N b\N q\N r/D).
\end{equation}
The coprimality with \(M\) is automatic on this support, since both
prime norms exceed \(\N M\).

Write \(B=\log z\), \(v(t)=\log(C/t)\), and
\[
 F_B(v)=\frac{2}{2B-v}\log\frac{B}{B-v}
       =\sum_{k\ge1}\alpha_k\frac{v^k}{B^{k+1}},
 \qquad \alpha_k=\sum_{j=1}^k\frac{2^{-(k-j)}}j>0.
\]
The quantitative fixed-field prime ideal theorem
\eqref{eq:quantitative-pnt}, applied twice with both upper cutoffs,
gives uniformly for \(b\mid M\)
\begin{equation}\label{eq:complete-cofactor-continuum}
 \sum_{\N q,\N r\le z}W(\N b\N q\N r/D)
 =\frac D{\N b}\int_c^C W(t)F_B(v(t)+\log\N b)\,dt
       +O_{\beta,W}\left(\frac D{\N b}
                       e^{-c_\beta\sqrt{\log D}}\right).
\end{equation}
Indeed, after setting \(xy=Dt/\N b\) and \(x=ze^{-s}\), the
continuous inner integral is
\[
 \int_0^{v(t)+\log\N b}
 \frac{ds}{(B-s)(B-v(t)-\log\N b+s)}
       =F_B(v(t)+\log\N b).
\]
For uniform errors, the first prime summation at \(x=\N q\)
costs \(O(D/(\N b\,x)e^{-c_\beta\sqrt{\log D}})\).
The sum of \(1/\N q\) costs at most a logarithm. The remaining
continuous function
\(G_b(x)=\int_2^z W(\N b\,xy/D)\,dy/\log y\) satisfies
\[
 |G_b(x)|\ll \frac D{\N b\,x\log D},\qquad
 |G_b'(x)|\ll \frac D{\N b\,x^2\log D}
\]
on its support. Partial summation proves the second error in the
same weighted form. Summing errors therefore costs
\(\sum_{b\mid M}1/\N b\), rather than \(\tau(M)\).
Removing the prime diagonal costs
\(O_W(\sqrt D\sum_{b\mid M}(\N b)^{-1/2})=O_\eps(D^{1/2+\eps})\).

Let \(M_j(W)=\int_c^C W(t)v(t)^j\,dt\), and choose the first
\(k_0\) with \(M_{k_0}(W)\ne0\). Such an index exists:
otherwise polynomial density in the compact logarithmic coordinate
would force \(W=0\). The zero integral gives \(k_0\ge1\).
Taylor expansion in bounded \(v\), followed by the vanishing lower
moments, gives uniformly for \(0\le l\le\beta\log D\)
\[
 \int W(t)F_B(v(t)+l)\,dt
 =\frac{\alpha_{k_0}M_{k_0}(W)}{B^{k_0+1}}
       +O_{\beta,W}\left(\frac{1+l}{B^{k_0+2}}\right).
\]
Here \(l/B\le2\beta+o(1)<1\), and
\(F_B^{(k_0+1)}(l+v)=O_{\beta,W}(B^{-k_0-2})\), which also
controls comparison of \(F_B^{(k_0)}(l)\) with its value at zero.

Put \(Z_+(M)=\prod_{p\mid M}(1+(\N p)^{-1})\) and
\(S(M)=\sum_{p\mid M}\log\N p/(\N p+1)\). Exact divisor
expansion gives
\[
 \sum_{b\mid M}\frac{\muK(b)}{\N b}=P_-(M),\quad
 \sum_{b\mid M}\frac1{\N b}=Z_+(M),\quad
 \sum_{b\mid M}\frac{\log\N b}{\N b}=Z_+(M)S(M).
\]
Consequently \eqref{eq:complete-cofactor-coherent} equals
\[
 \frac{D\alpha_{k_0}M_{k_0}(W)P_-(M)}{B^{k_0+1}}
 +O_{\beta,W}\left(\frac{DZ_+(M)(1+S(M))}{B^{k_0+2}}
       +DZ_+(M)e^{-c_\beta\sqrt{\log D}}\right)
 +O_\eps(D^{1/2+\eps}).
\]
To control this uniformly, split the prime divisors of \(M\) at
\(T=\log D\). The fixed-field prime ideal theorem gives
\(\sum_{\N p\le T}1/\N p=\log\log T+O(1)\) and
\(\sum_{\N p\le T}\log\N p/\N p=\log T+O(1)\).
Above \(T\), use \(\sum_{p\mid M}\log\N p\le\beta\log D\).
It follows that
\[
 P_-(M)\gg_\beta(\log\log D)^{-1},\qquad
 Z_+(M)/P_-(M)\ll_\beta(\log\log D)^2,\qquad
 S(M)\ll_\beta\log\log D.
\]
The relative error is \(O_{\beta,W}((\log\log D)^3/\log D)=o(1)\).
Thus \(|B_{v^6,M}|\gg_{\beta,W}DP_-(M)/(\log D)^{k_0+1}\).

Finally, uniform lattice counting in each principal ideal and divisor
inclusion-exclusion give, for \(X=D^{1/15}\),
\[
 \#\{0<\N v\le X:(v,M)=1\}
 =\kappa_K XP_-(M)+O_K(\tau(M)(\sqrt X+1)).
\]
Use \(\tau(M)\ll_\eps D^\eps\), with \(\eps<1/30\), to
make the error negligible uniformly. The sixth-power map has at most
six preimages. Squaring on these rows, using \(\Phi\ge1\) on the
inner ball, and dividing by \(D\) proves the theorem.
\end{proof}

The essential Jacobian is \(1/\N b\): its signed divisor sum is
\(P_-(M)>0\), although \(\sum_{b\mid M}\muK(b)=0\) for
nonunit \(M\). Complete cofactor pairing here therefore retains the
power \(16/15\), even for every fixed nonzero zero-integral profile.
Writing the full response as \(T_{u,\mathrm{sf}}=B_{u,M}+R_{u,M}\)
leaves an uncontrolled cross term. Balanced triples with all prime
norms near \(D^{1/3}\) lie outside this block and can still contribute
to its compensation. No full-tail lower bound follows.

\subsection{Grouping the Möbius factors does not repair the generic bound}

Group \(d=ab\), and take a dyadic rectangle
\(\N d\asymp X\), \(\N m\asymp Z\), \(XZ\asymp D\).
The coefficients \(c_z(d)\) are cube-free-supported, hence
sixth-power-free, with squared norm \(O_\eps(XD^\eps)\).
Cauchy in \(m\) and the physical generic sieve
\eqref{eq:physical-operator} gives only
\begin{equation}\label{eq:generic-bilinear}
 \mathfrak E_{X,Z}\ll_\eps D^\eps\frac D X
 \left[H+XH^{1/6}+H^{5/6}X^{1/3}+H^{1/3}X^{5/6}\right].
\end{equation}
Indeed Cauchy pays \(Z\), summing the \(m\)-indexed sieve bounds
pays another \(Z\), the \(d\)-coefficient norm pays \(X\), and
normalization by \(D\asymp XZ\) leaves \(Z\) times the operator factor.
For each fixed \(m\), the dyadic profile has bounded seminorms.

The threshold \(\N d>Y_u\) in this particular grouping is covered
by a maximal partial-sum form of the sieve with only a
\(\log^2(2X)\) loss. Order the columns by norm with fixed tie-breaking.
Every suffix decomposes into \(O(\log(2X))\) intervals from a fixed
binary partition. Cauchy and the sieve on each fixed interval,
summed by partition level, give the asserted loss because each level
partitions the coefficient squared norm.
This does not authorize arbitrary selectors in completed row kernels.

For \(H=D^h\), \(X=D^x\), the four supplied powers are
\[
 1+h-x,\qquad 1+\frac h6,\qquad
 1+\frac{5h}{6}-\frac{2x}{3},\qquad
 1+\frac h3-\frac x6.
\]
The term \((D/X)XH^{1/6}=DH^{1/6}\) survives every balance.
This is an insufficiency of the generic upper bound, not a lower
bound for the actual M\"obius block.
The \(m=1\) block, \(X\asymp D\), remains in the tail.
Its coefficient norm has the full power \(D\): on any subinterval
where \(|W|\) is bounded below, the prime ideal theorem gives
\(\gg D/\log^2D\) balanced distinct semiprimes \(pq\) with
\(c_z(pq)=2\).
Zero integral provides no power saving in that norm.

\section{Routes without an additional power gain and the next task}
\label{sec:limits}

Several natural alterations simplify a presentation without supplying
the arithmetic saving in \eqref{eq:intro-moment}.
The replication identities and effective row-count bound following
Theorem~\ref{thm:finite-extraction} show why composite sixth powers,
free row weights, and lower-order subfamilies inherited solely from
the same moment keep the same extraction floor. Their full
bookkeeping is in \cite{N5}.
The core-cutoff estimates above quantify the extra conductor
uniformity needed when the collection of fixed cores grows.

The source's positive dual estimate cannot simply be extended into
the long-row regime. At column length \(X=1\), a smooth profile
isolating the unit ideal gives a positive dual energy
\(\mathcal E(R,1,F)\asymp R\), with the actual source coefficient
\(a_\xi(1)=1\).
For \(b=1,F=D,H=D^h\), it has \(R\asymp D^{2-h}\),
whereas the proposed old right side has scale \(\Sigma=XF=D\).
This obstructs that positive extension for \(h<1\).
The full diagonal removal avoids this particular unit term;
it does not prove the remaining signed estimate.
See \cite{N2} for the exact test and \cite{N1} for the failed
shrinking-ratio condition.

A second Poisson transform of the coprime core is the exact involution
proved above, and unit-character splitting has only finitely many
classes \cite{N3}. Neither gives a power contraction.
The positive divisor moments in the gcd identity retain
\(B_1=\Mfrak_W\); separately triangle-bounding them is circular.
Finally, a Poisson sum over either remaining M\"obius factor would
require information about its coefficients that the free-factor
completion does not supply.

One concrete sufficient input remains the following.

\begin{assumption}[Adaptive signed tail moment]\label{input:open-tail}
For fixed \(h=4/5\), some \(0\le a<1/12\), and one fixed
\(0<\theta<1\), every \(V\in C_c^\infty((0,\infty))\) with
\(W=(1+t\partial_t)V\) satisfies, for every \(\eps>0\),
\[
 \frac1D\sum_{u\ne0}\Phi(\N u/D^{4/5})
 \left|\sum_{\N d>D^{1-\theta}/L_u}
     c_z(d)\lambda_u(d)S_u(D/\N d)\right|^2
       \ll_{\eps,V}D^{4/5+a+\eps}.
\]
All \(z\) cutoffs, overlapping factors, deletion zeros, and the
complete row weight are retained.
\end{assumption}

Theorem~\ref{sfac:thm:adaptive} and
Theorem~\ref{thm:finite-extraction} would turn this input into
scalar exponent \(5/6+a/2<7/8\).
The derivative detectors then exclude zeros to its right,
with \(s=1\) treated separately.
The alternative \(h=8/9\), \(a<1/108\), gives
\(47/54+a/2<7/8\).
At \(h=a=2/5\), the proposed moment would give \(13/15<7/8\).
Corollary~\ref{cor:sf-residual} reduces that target to its squarefree
signed tail at \(\theta=11/20\). None of these full moments is
established.

For an endpoint program, moments at a sequence
\(h_j\to0\), \(a_j\to0\), for every detecting profile, would give
scalar exponents approaching \(1/2\).
In fact their \(u=1\) terms alone already approach that endpoint;
finite replication sharpens the finite-\(h\) output.
For \(\nu=1\), zero-freeness of
\(\zeta_K(s)=\zeta(s)L(s,\chi_{-3})\) in \(\Rea s>1/2\)
would imply RH for \(\zeta\) by its functional equation.
A zeta-only quasi-RH hypothesis supplies none of the new
Hecke-family moments automatically.

\appendix
\section{Research record and scope of verification}\label{sec:record}

This manuscript organizes the working-tree results developed through
8 October 2026, based on repository commit
\texttt{480581447dcd3e93c9c2f9c22a050d9c5100ff9d}.
The notes contain fuller source audits, conventions and exploratory
calculations. They are grouped here by their role.

\begin{center}\small
\begin{tabular}{@{}p{.14\textwidth}p{.76\textwidth}@{}}
\toprule
Notes & Role in the organized argument\\
\midrule
1 & Finite reuse of one moment, prime-mask extraction and Mellin detection.\\
2--4 & Full diagonal removal, exact involution, complete primitive kernels
and the enlarged overlap sector.\\
5--7 & Replication limits, generic sextic-sieve transfer and the earlier
signed residual.\\
8--10 & Exact truncated inverse, centered uniform cutoff, signed gcd
truncation and quantitative coherent-core costs.\\
11--13 & Zero-integral profiles, sharp-profile recovery, radical-dependent
adaptive reduction, factor-selected obstruction and the current target.\\
15 & Exact nonunit-cofactor packet cancellation, a growing full-weight
controlled sector, opposite branch lower bounds and the surviving
semiprime coefficient.\\
17 & Whole-nonsquarefree sector bound, sparse adaptive-prefix estimate
and the equivalent squarefree residual at \(h=a=2/5\).\\
21 & Complete small-prime cofactor parity, a growing odd-core zero sector
at large conductor, and the unchanged retained-response budget.\\
24 & A nonzero complete divisor-cofactor pairing retains the coherent
power uniformly for cofactor norm at most \(D^\beta\), \(\beta<1/20\).\\
\bottomrule
\end{tabular}
\end{center}

The existing finite checkers use exact rational or finite cyclotomic
arithmetic. The consolidated records include 9,356 assertions for
the core/overlap continuation, 24,245 for factorization and gcd
algebra, and 4,812 for the zero-integral adaptive continuation.
The divisor-packet record adds 3,982 exact assertions, including
nonsquarefree cofactors, cutoff boundaries, complex masks and actual
local split-prime sextic symbols. Its local cutoff uses a good-radical
proxy, not a computation of the inducing primitive conductor, and checks
the pointwise packet theorem rather than its asymptotic growing range.
The squarefree-projection record adds 594,360 exact assertions for
projection signs, equal-norm distinct ideal symbols, square-divisor
overlaps, sixth-power decomposition and fixed binary prefixes.
These finite checks do not certify the asymptotic sparse count or
operator estimate used in the sector theorem.
The rough-parity record adds 32,815 assertions comparing original ordered
factor tuples with complete cofactors, saturation, overflow, cutoff
endpoints and nonempty signed zero sectors. These do not certify an
analytic moment estimate.
The complete-cofactor record adds 3,289 assertions for actual ordered
tuple coefficients, phase and deletion zeros, dilation Euler products,
Jacobian factors, the balanced-triple remainder and rational margins.
It does not certify the uniform prime ideal or lattice asymptotics
in Theorem~\ref{thm:complete-cofactor-obstruction}.
They check specified finite identities, complex conjugation
orientations, deletion zeros and rational exponent arithmetic.
They do not verify the imported Poisson theorem, infinite Euler
convergence, quantitative prime ideal theorem, generic sieve,
or an asymptotic signed moment. Detailed review links are in
the notes and repository index.

\subsection*{Acknowledgement and preparation metadata}

Prepared for Edward Baker with substantial assistance from GPT-6
(Codex) in organization, derivation, source comparison and internal
review. The inherited configuration does not expose the exact serving
variant or configured reasoning effort, and neither is inferred.
Parallel same-model audits and finite checks are internal validation,
not independent specialist review or formal proof verification.
Authorship and any publication version require the author's review.

\begin{thebibliography}{OA30}
\bibitem[OA5]{OA5}
OpenAI, \emph{The Quasi-Riemann Hypothesis}, version dated
5 October 2026, source archive.
\href{https://github.com/openai/math/tree/main/preprints/The-Quasi-Riemann-Hypothesis-October-5-2026}{Online source}.
Native character and Poisson formulas and the positive dual estimate
are used only in the explicitly recorded scope.
\bibitem[DF]{DF}
Alexandre de Faveri, \emph{Optimal large sieve for fixed order characters},
arXiv:2610.04045v1, 2 October 2026.
\href{https://arxiv.org/html/2610.04045v1}{Version 1}.
\bibitem[DKN]{DKN}
Sourabhashis Das, Habiba Kadiri and Nathan Ng,
\emph{An effective version of Chebotarev's density theorem},
arXiv:2508.09480v1, 13 August 2025.
\href{https://arxiv.org/html/2508.09480v1}{Version 1}.

\bibitem[N1]{N1}
Research note 1, \emph{A fixed short-family moment already removes its prime-mask error}, 8 October 2026.
\href{notes/1_SHORT_FAMILY_DESCENT_REFINEMENT_20261008.md}{Repository note 1}.
\bibitem[N2]{N2}
Research note 2, \emph{Removing the diagonal and the large transformed cofactor}, 8 October 2026.
\href{notes/2_SHORT_FAMILY_SMALL_COFACTOR_20261008.md}{Repository note 2}.
\bibitem[N3]{N3}
Research note 3, \emph{Source corrections, exact involution, and a prime-column obstruction}, 8 October 2026.
\href{notes/3_SHORT_FAMILY_CORE_INVOLUTION_20261008.md}{Repository note 3}.
\bibitem[N4]{N4}
Research note 4, \emph{A uniform primitive-kernel bound and a larger controlled overlap sector}, 8 October 2026.
\href{notes/4_SHORT_FAMILY_HIGH_OVERLAP_20261008.md}{Repository note 4}.
\bibitem[N5]{N5}
Research note 5, \emph{Limits of composite replication, row weights, and inherited lower-order families}, 8 October 2026.
\href{notes/5_SHORT_FAMILY_REPLICATION_LIMITS_20261008.md}{Repository note 5}.
\bibitem[N6]{N6}
Research note 6, \emph{The generic sixth-order sieve reaches the coherent-row barrier}, 8 October 2026.
\href{notes/6_SHORT_FAMILY_SEXTIC_SIEVE_BARRIER_20261008.md}{Repository note 6}.
\bibitem[N7]{N7}
Research note 7, \emph{Short families: the remaining cancellation problem}, 8 October 2026.
\href{notes/7_SHORT_FAMILY_CONTINUATION_20261008.md}{Repository note 7}.
\bibitem[N8]{N8}
Research note 8, \emph{Exact M\"obius factorization and a controlled centered small product}, 8 October 2026.
\href{notes/8_SHORT_FAMILY_CENTERED_FACTORIZATION_20261008.md}{Repository note 8}.
\bibitem[N9]{N9}
Research note 9, \emph{Signed gcd recombination of the original coprime core}, 8 October 2026.
\href{notes/9_SHORT_FAMILY_SIGNED_GCD_RECOMBINATION_20261008.md}{Repository note 9}.
\bibitem[N10]{N10}
Research note 10, \emph{Continuation after centered factorization}, 8 October 2026.
\href{notes/10_SHORT_FAMILY_FACTORIZATION_CONTINUATION_20261008.md}{Repository note 10}.
\bibitem[N11]{N11}
Research note 11, \emph{Zero-integral profiles and a conductor-dependent factorization cutoff}, 8 October 2026.
\href{notes/11_SHORT_FAMILY_MEAN_ZERO_ADAPTIVE_REDUCTION_20261008.md}{Repository note 11}.
\bibitem[N12]{N12}
Research note 12, \emph{Factor selection on zero-integral profiles}, 8 October 2026.
\href{notes/12_SHORT_FAMILY_FACTOR_SELECTION_BARRIER_20261008.md}{Repository note 12}.
\bibitem[N13]{N13}
Research note 13, \emph{Continuation after the zero-integral adaptive reduction}, 8 October 2026.
\href{notes/13_SHORT_FAMILY_ADAPTIVE_CONTINUATION_20261008.md}{Repository note 13}.
\bibitem[N15]{N15}
Research note 15, \emph{Exact divisor packet cancellation in the adaptive tail}, 8 October 2026.
\href{notes/15_SHORT_FAMILY_DIVISOR_PACKET_CANCELLATION_20261008.md}{Repository note 15}.
\bibitem[N17]{N17}
Research note 17, \emph{A power bound for the nonsquarefree tail and a squarefree residual}, 8 October 2026.
\href{notes/17_SHORT_FAMILY_SQUAREFREE_TAIL_REDUCTION_20261008.md}{Repository note 17}.
\bibitem[N21]{N21}
Research note 21, \emph{Squarefree cofactor parity and a growing signed zero sector}, 8 October 2026.
\href{notes/21_SHORT_FAMILY_ROUGH_PARITY_20261008.md}{Repository note 21}.
\bibitem[N24]{N24}
Research note 24, \emph{A complete small-cofactor pairing that retains the coherent power}, 8 October 2026.
\href{notes/24_SHORT_FAMILY_FINITE_COFACTOR_PAIRING_20261008.md}{Repository note 24}.
\end{thebibliography}
\end{document}
